%Version 3.1 December 2024
% See section 11 of the User Manual for version history
%
%%%%%%%%%%%%%%%%%%%%%%%%%%%%%%%%%%%%%%%%%%%%%%%%%%%%%%%%%%%%%%%%%%%%%%
%%                                                                 %%
%% Please do not use \input{...} to include other tex files.       %%
%% Submit your LaTeX manuscript as one .tex document.              %%
%%                                                                 %%
%% All additional figures and files should be attached             %%
%% separately and not embedded in the \TeX\ document itself.       %%
%%                                                                 %%
%%%%%%%%%%%%%%%%%%%%%%%%%%%%%%%%%%%%%%%%%%%%%%%%%%%%%%%%%%%%%%%%%%%%%

%%\documentclass[referee,sn-basic]{sn-jnl}% referee option is meant for double line spacing

%%=======================================================%%
%% to print line numbers in the margin use lineno option %%
%%=======================================================%%

%%\documentclass[lineno,pdflatex,sn-basic]{sn-jnl}% Basic Springer Nature Reference Style/Chemistry Reference Style

%%=========================================================================================%%
%% the documentclass is set to pdflatex as default. You can delete it if not appropriate.  %%
%%=========================================================================================%%

%%\documentclass[sn-basic]{sn-jnl}% Basic Springer Nature Reference Style/Chemistry Reference Style

%%Note: the following reference styles support Namedate and Numbered referencing. By default the style follows the most common style. To switch between the options you can add or remove "Numbered" in the optional parenthesis. 
%%The option is available for: sn-basic.bst, sn-chicago.bst%  
 
%%\documentclass[pdflatex,sn-nature]{sn-jnl}% Style for submissions to Nature Portfolio journals
%%\documentclass[pdflatex,sn-basic]{sn-jnl}% Basic Springer Nature Reference Style/Chemistry Reference Style
\documentclass[pdflatex,sn-basic,Numbered,referee,lineno]{sn-jnl}% J Cheminformatics: chemistry ref style, numbered, double-spaced, line numbers
%%\documentclass[pdflatex,sn-mathphys-ay]{sn-jnl}% Math and Physical Sciences Author Year Reference Style
%%\documentclass[pdflatex,sn-aps]{sn-jnl}% American Physical Society (APS) Reference Style
%%\documentclass[pdflatex,sn-vancouver-num]{sn-jnl}% Vancouver Numbered Reference Style
%%\documentclass[pdflatex,sn-vancouver-ay]{sn-jnl}% Vancouver Author Year Reference Style
%%\documentclass[pdflatex,sn-apa]{sn-jnl}% APA Reference Style
%%\documentclass[pdflatex,sn-chicago]{sn-jnl}% Chicago-based Humanities Reference Style

%%%% Standard Packages
%%<additional latex packages if required can be included here>

\usepackage{graphicx}%
\usepackage{multirow}%
\usepackage{amsmath,amssymb,amsfonts}%
\usepackage{amsthm}%
\usepackage{mathrsfs}%
\usepackage[title]{appendix}%
\usepackage{xcolor}%
\usepackage{textcomp}%
\usepackage{manyfoot}%
\usepackage{booktabs}%
\usepackage{algorithm}%
\usepackage{algorithmicx}%
\usepackage{algpseudocode}%
\usepackage{listings}%
%%%%

%%%%%=============================================================================%%%%
%%%%  Remarks: This template is provided to aid authors with the preparation
%%%%  of original research articles intended for submission to journals published 
%%%%  by Springer Nature. The guidance has been prepared in partnership with 
%%%%  production teams to conform to Springer Nature technical requirements. 
%%%%  Editorial and presentation requirements differ among journal portfolios and 
%%%%  research disciplines. You may find sections in this template are irrelevant 
%%%%  to your work and are empowered to omit any such section if allowed by the 
%%%%  journal you intend to submit to. The submission guidelines and policies 
%%%%  of the journal take precedence. A detailed User Manual is available in the 
%%%%  template package for technical guidance.
%%%%%=============================================================================%%%%

%% as per the requirement new theorem styles can be included as shown below
\theoremstyle{thmstyleone}%
\newtheorem{theorem}{Theorem}%  meant for continuous numbers
%%\newtheorem{theorem}{Theorem}[section]% meant for sectionwise numbers
%% optional argument [theorem] produces theorem numbering sequence instead of independent numbers for Proposition
\newtheorem{proposition}[theorem]{Proposition}% 
%%\newtheorem{proposition}{Proposition}% to get separate numbers for theorem and proposition etc.

\theoremstyle{thmstyletwo}%
\newtheorem{example}{Example}%
\newtheorem{remark}{Remark}%

\theoremstyle{thmstylethree}%
\newtheorem{definition}{Definition}%

\raggedbottom
\unnumbered% unnumbered section headings -- BMC / J Cheminformatics house style

\begin{document}



%% [TODO] 대괄호는 러닝 헤드(짧은 제목)다.
\title[Calibration collapse in time-split validation]{Beyond AUROC: Quantifying
Calibration Collapse in Time-Split Validation of Hepatotoxicity and
Cardiotoxicity Models}

%%=============================================================%%
%% GivenName	-> \fnm{Joergen W.}
%% Particle	-> \spfx{van der} -> surname prefix
%% FamilyName	-> \sur{Ploeg}
%% Suffix	-> \sfx{IV}
%% \author*[1,2]{\fnm{Joergen W.} \spfx{van der} \sur{Ploeg} 
%%  \sfx{IV}}\email{iauthor@gmail.com}
%%=============================================================%%

%% 단독 저자. 템플릿의 둘째·셋째 저자와 \equalcont 는 지웠다.
%% [TODO] 이름·이메일을 채울 것. ORCID 는 투고 시스템에 입력한다.
\author*[1]{\fnm{[이름]} \sur{[성]}}\email{[이메일]}

%% [TODO] 소속 한 곳.
\affil*[1]{\orgdiv{[학과]}, \orgname{[기관]}, \orgaddress{\city{[도시]}, \country{[국가]}}}

%%==================================%%
%% Sample for unstructured abstract %%
%%==================================%%

%%===========================================================%%
%% Journal of Cheminformatics: abstract must be UNSTRUCTURED, %%
%% 350 words maximum, and must contain a section titled      %%
%% "Scientific Contribution" of at most THREE sentences.     %%
%%===========================================================%%

\abstract{[Aim, approach and principal findings, in continuous prose. No
subheadings, no citations, no equations. 350 words maximum including the
sentences below.]

\textbf{Scientific Contribution} [Sentence 1: what this work establishes that
was not established before.] [Sentence 2: how it differs from prior work on
this topic.] [Sentence 3: why that difference matters for how these models are
validated. --- Three sentences maximum.]}

%%================================%%
%% Sample for structured abstract %%
%%================================%%

% \abstract{\textbf{Purpose:} The abstract serves both as a general introduction to the topic and as a brief, non-technical summary of the main results and their implications. The abstract must not include subheadings (unless expressly permitted in the journal's Instructions to Authors), equations or citations. As a guide the abstract should not exceed 200 words. Most journals do not set a hard limit however authors are advised to check the author instructions for the journal they are submitting to.
% 
% \textbf{Methods:} The abstract serves both as a general introduction to the topic and as a brief, non-technical summary of the main results and their implications. The abstract must not include subheadings (unless expressly permitted in the journal's Instructions to Authors), equations or citations. As a guide the abstract should not exceed 200 words. Most journals do not set a hard limit however authors are advised to check the author instructions for the journal they are submitting to.
% 
% \textbf{Results:} The abstract serves both as a general introduction to the topic and as a brief, non-technical summary of the main results and their implications. The abstract must not include subheadings (unless expressly permitted in the journal's Instructions to Authors), equations or citations. As a guide the abstract should not exceed 200 words. Most journals do not set a hard limit however authors are advised to check the author instructions for the journal they are submitting to.
% 
% \textbf{Conclusion:} The abstract serves both as a general introduction to the topic and as a brief, non-technical summary of the main results and their implications. The abstract must not include subheadings (unless expressly permitted in the journal's Instructions to Authors), equations or citations. As a guide the abstract should not exceed 200 words. Most journals do not set a hard limit however authors are advised to check the author instructions for the journal they are submitting to.}

%% [TODO] 저널 지침 확인 후 확정.
\keywords{calibration, temporal validation, drug-induced liver injury,
drug-induced cardiotoxicity, conformal prediction, model evaluation}

%%\pacs[JEL Classification]{D8, H51}

%%\pacs[MSC Classification]{35A01, 65L10, 65L12, 65L20, 65L70}

\maketitle


%====================================================================%
% 조립 순서: Introduction - Results - Discussion/Conclusions - Methods %
% 각 부분은 별도 파일이다. **이 파일에 산문을 직접 쓰지 말 것** —      %
% 부분 파일을 고치고 여기서는 \input 만 유지한다.                      %
%                                                                    %
%   introduction.tex      아직 채울 자리 (산문 없음)                   %
%   results_body_v2.tex   현행 (verify12 통과)                        %
%   tables_main.tex       src/make_tables.py 가 생성                  %
%   figures.tex           범례 (PDF 는 src/make_figures.py)            %
%   discussion_body_v1.tex 현행 — 신규 절 반영 필요                    %
%   methods_body.tex      현행 (2026-10-20 표본수·반복수 수정)         %
%   declarations.tex      저널 지침 6항목                              %
%====================================================================%

%====================================================================%
% Introduction                                                       %
%                                                                    %
% 아직 산문이 없다. 아래는 **채울 자리**이고, 각 단락에 무엇을 쓸지와   %
% 어느 근거를 인용할지를 적어 두었다. 수치를 쓸 때는 보고서            %
% (실험결과_논문용_*.md)의 해당 절을 보고 그 값을 그대로 옮긴 뒤       %
% `python src/verify12.py` 로 대조할 것.                              %
%====================================================================%

\section{Introduction}

% [단락 1 — 문제] 독성 예측 모델은 보통 무작위 분할로 평가되고, 보고되는
% 지표는 거의 언제나 AUROC 다. 그런데 실제 사용은 전향적이다 — 과거 약물로
% 학습해 새 약물에 쓴다. 이 괴리를 한 문장으로 세울 것.

% [단락 2 — 알려진 것] 시간분할이 무작위 분할보다 어렵다는 것은 이미
% 알려져 있다. 선행연구를 인용하고, **그 문헌이 무엇을 보고했는지**(대개
% 판별력 저하)를 적을 것. 참고문헌은 02_참고문헌_목록.md 에 정리돼 있다.

% [단락 3 — 빈 자리] 판별력 저하만 보고하면 놓치는 것이 있다. 확률의
% 크기를 믿고 쓰는 쪽(우선순위 결정, 임계값 운용)에서는 보정이 더
% 중요한데, 보정 붕괴를 분리해 재고 그것이 유병률 변화로 설명되는지까지
% 통제한 연구가 없다는 점을 적을 것.

% [단락 4 — 이 논문이 하는 것] 두 데이터셋(DILIrank, DICTrank)에서
% 2010년 시간분할을 세우고, (i) Brier 분해로 판별력과 보정을 분리하고,
% (ii) 유병률을 맞춘 무작위 대조군 R1--R3 로 유병률 효과를 걷어내고,
% (iii) 재보정·conformal·최신성 가중 등 처방을 전향적 제약 아래에서
% 시험했다는 것을 적을 것. 숫자는 Results 로 넘기고 여기서는 설계만.

% [단락 5 — 결과 한 줄 예고] 저하의 대부분이 보정이고, 유병률을 양쪽
% 통제해도 남으며, 사후 재보정으로는 일부만 되돌아온다는 것. 구체 수치는
% 쓰지 말 것(초록과 Results 가 담당한다).



%%===================================================================%%
%% Results — revision of 2026-09-21                                  %%
%% Supersedes results_body.tex (2026-09-19).                         %%
%% Every number here is traceable to a file under results/ produced  %%
%% by rerun_all.sh on the 796 / 919 analysis sets.                   %%
%%===================================================================%%

\section{Results}

\subsection{Temporal degradation in both datasets}

The analysis sets comprise 796 DILIrank and 919 DICTrank compounds with an
approval year, a standardised structure and a binary label. Splitting at 2010
leaves 640 training and 156 held-out compounds for DILIrank, and 652 and 267 for
DICTrank.

On the DILIrank held-out set the model reached an AUROC of 0.648 (95\% bootstrap
CI 0.555--0.736). One thousand random partitions of the same held-out size
averaged 0.735 (SD 0.036, 2.5--97.5 percentile 0.663--0.806), and six of them fell
at or below the observed value. Using the bias-corrected estimator of Phipson and
Smyth, which cannot return zero, the one-sided permutation $p$ is 0.0070.
DICTrank behaved the same way: 0.620 (0.553--0.688) against a random reference of
0.697 (SD 0.029, 0.636--0.752), four partitions at or below, $p = 0.0050$. Both
survive Benjamini--Hochberg correction within the confirmatory family at
$q = 0.0070$.

The held-out sets are also less positive than the training sets, so part of the
AUROC gap could be an artefact of class balance. Resampling the DILIrank held-out
set to the training positive rate of 65.5\%, averaged over 20 draws, moved AUROC
from 0.648 to 0.632 on the resulting 73 compounds --- that is, downwards. Measured
against the 20-partition random control used in the same script, which gives
0.742, the gap widens from 0.094 to 0.110 rather than closing. The direction
matters: class
balance is not masking a smaller gap, and the shift accounts for none of the
degradation in discrimination. Its effect on calibration is a separate question
and is treated below.

Gaps of 0.09 and 0.08 are modest. Reported alone they would say only that
prospective validation is harder than retrospective validation, which is already
well known. The remainder of this section concerns what those numbers conceal.

\subsection{The degradation is a calibration loss}

Table~\ref{tab:brier} gives the Murphy decomposition. Each random partition is
decomposed at the same held-out size as the temporal split and the terms are then
averaged, so the two columns rest on equal sample sizes.

For DILIrank the Brier score rose from 0.213 to 0.427. Resolution, which tracks
discrimination, fell by 0.034. Reliability, which tracks calibration, rose by
0.208. Calibration therefore accounts for \textbf{86.1\%} (95\% CI 84.6--97.3) of
the degradation against 13.9\% for discrimination, a ratio of 6.2 to 1. DICTrank
gives \textbf{82.5\%} (75.5--124.0) against 17.5\%, a ratio of 4.7 to 1: a Brier
score of 0.205 rising to 0.306, resolution down 0.012, reliability up 0.054.

Pooling the random partitions before binning, as an earlier version of this
analysis did, inflates both shares --- to 90.4\% for DILIrank and 91.8\% for
DICTrank. Pooling gives the random reference twenty times the sample of the
temporal column, which shrinks its reliability term and flatters the comparison by
4.3 and 9.3 percentage points. The per-partition figures are reported as primary
throughout.

The ordinary calibration metrics agree. Against a random split of the same size,
expected calibration error rose from 0.121 to 0.438 in DILIrank and from 0.126 to
0.270 in DICTrank, and maximum calibration error from 0.295 to 0.622 and from
0.471 to 0.593. The calibration slope fell from 0.534 to 0.270 and from 0.433 to
0.222, and the intercept moved from $+0.053$ to $-1.313$ and from $+0.326$ to
$-0.085$. Mean predicted probability on the DILIrank held-out set was 0.732 where
the observed positive rate was 0.308. For DICTrank the figures are 0.845 and
0.614.

These are over-confident predictions, and the over-confidence is worst in the bins
a user would act on. Seventy-one of the 156 DILIrank compounds were assigned a
probability above 0.9, with a mean prediction of 0.962; 40.8\% of them were
positive. The corresponding DICTrank bin held 156 of 267 compounds at a mean
prediction of 0.966 against an observed 67.9\%. Figure~\ref{fig:calib} shows the
reliability curves in full.

Because the failure is a shift rather than a scrambling, moving the decision
threshold recovers part of what is lost at 0.5. The threshold maximising Youden's
J was 0.920 for DILIrank and 0.940 for DICTrank. At those thresholds balanced
accuracy rose from 0.557 to 0.635 and from 0.519 to 0.601, and the Matthews
correlation coefficient from 0.126 to 0.254 and from 0.062 to 0.198. A threshold
near 0.93 is not a usable default; the point is that one has to move that far to
find it.

AUROC is invariant to any monotone transformation of the predicted probabilities,
so none of this appears in it.

\subsection{Scaffold splitting does not reproduce the calibration collapse}

Compounds approved after 2010 are structurally newer as well as later, which
raises the possibility that any sufficiently hard split would produce the same
picture. Running stratified and Bemis--Murcko scaffold cross-validation on the
same compounds, model and features tests that directly
(Table~\ref{tab:split}, Figure~\ref{fig:split}a,b).

Discrimination gave an inconsistent answer. Scaffold cross-validation raised
DILIrank AUROC slightly, from 0.728 $\pm$ 0.022 to 0.731 $\pm$ 0.018, and lowered
DICTrank from 0.699 $\pm$ 0.023 to 0.678 $\pm$ 0.036. Both changes are smaller
than the fold-to-fold standard deviation, and they point in opposite directions in
the two datasets. Read through AUROC alone, scaffold splitting is neither reliably
harder nor reliably easier than a random split here.

Calibration gave a consistent one. Scaffold ECE was 0.130 $\pm$ 0.029 against
0.136 $\pm$ 0.012 for DILIrank and 0.136 $\pm$ 0.030 against 0.129 $\pm$ 0.022 for
DICTrank --- changes of $-0.006$ and $+0.007$, within a standard deviation in both
cases. The temporal split on the same compounds gave 0.438 and 0.270, which is
3.4 times the scaffold value in DILIrank and 2.0 times it in DICTrank. Structural
novelty in the held-out set, imposed deliberately, leaves calibration where a
random split leaves it; temporal ordering does not.

One reason the two cross-validation designs resemble each other is visible in the
scaffold counts. The 796 DILIrank compounds contain 542 unique Bemis--Murcko
scaffolds, 455 of which appear exactly once; the 919 DICTrank compounds contain
638, with 545 singletons. When 84\% and 85\% of scaffolds are singletons,
grouping by scaffold removes little that a random partition would not already
separate, and the design converges on stratified cross-validation. This is a
property of drug-like reference sets rather than of these two datasets in
particular, and it is worth stating plainly before scaffold splitting is offered
as a proxy for prospective validation.

This comparison was not among the pre-specified hypotheses, and it is reported as
a description rather than a test. No $p$-value is attached to it.

\subsection{Prevalence-matched controls}

The positive rate falls across the split boundary, from 65.5\% to 30.8\% in
DILIrank and from 77.0\% to 61.4\% in DICTrank, and a model calibrated to one base
rate will be miscalibrated on another. Three nested controls, each averaged over
200 partitions, separate that effect from the rest (Figure~\ref{fig:split}c).
Because they use a different number of repeats from the calibration analysis
above, their random-split values differ from it in the second decimal place.

DILIrank temporal ECE was 0.438 (95\% CI 0.374--0.527). An unconstrained random
split (R1) gave 0.131 $\pm$ 0.029, a ratio of 3.35. Matching the held-out positive
rate (R2) raised the random-split value to 0.250 $\pm$ 0.032, ratio 1.76. Matching
both rates (R3) raised it to 0.321 $\pm$ 0.032, ratio 1.37. DICTrank temporal ECE
was 0.270 (0.228--0.336), against 0.130, 0.187 and 0.223 for ratios of 2.08, 1.45
and 1.21.

No random partition out of 200 exceeded the temporal ECE in five of the six
comparisons, giving $p = 0.0050$, the lower bound available at 200 repeats under
the Phipson--Smyth estimator. The exception is DICTrank R3, where six of 200 did,
giving $p = 0.0348$. All six survive correction within their family, the largest
$q$ being 0.0348.

Matching both base rates closes most but not all of the excess: of the 0.308 of
ECE separating the temporal split from an unconstrained random one in DILIrank,
0.117 survives R3, a closure of 62\%; of the 0.140 in DICTrank, 0.048 survives, a
closure of 66\%. Prevalence shift is a large part of the calibration collapse and
not the whole of it, and the two datasets agree closely on how large a part.

Discrimination behaves differently under the same controls. R3 AUROC was 0.726 for
DILIrank and 0.693 for DICTrank against temporal values of 0.648 and 0.620, so
matching the base rates removes little of the AUROC gap while removing most of the
ECE gap. In DILIrank the fully matched comparison is a 10.8\% relative loss in
AUROC against a 36.5\% relative rise in ECE.

\subsection{Recalibration under a prospective constraint}

Holding out the most recent quarter of the training compounds, ordered by approval
year, as an internal calibration set keeps strategies (B) and (C) inside the
information a prospective user would have; strategy (D) does not, for the reason
given below. Table~\ref{tab:recal} gives the result, and it is a mixed one.

DILIrank began at ECE 0.447. The oracle ceiling, Platt scaling fitted on the
held-out labels, reached 0.048. Internal Platt scaling closed 8\% of that gap and
isotonic regression 10\%. Only the prior-probability correction moved the needle,
reaching 0.296 and closing 38\%. DICTrank reversed the ordering: internal Platt
scaling reached 0.168 and closed 42\%, isotonic regression 0.174 and 40\%, and the
prior correction 0.244 and 12\% against an oracle of 0.023.

The prior-probability correction is not fully prospective and is reported with
that caveat. It rescales the intercept using the held-out positive rate, which a
user would have to estimate rather than know. It is included because it isolates
how much of the miscalibration is a pure base-rate effect; strategies (B) and (C)
are the ones that use no held-out information at all.

Neither of the two fully prospective strategies recovered even half of the
recoverable error, and the two datasets differ by a factor of four to five in how
much they recover at all: 8--10\% for DILIrank against 40--42\% for DICTrank. The
base-rate correction runs the other way, recovering 38\% in DILIrank and 12\% in
DICTrank. The positive rate moves by 35 percentage points across the DILIrank
split boundary and by 16 across the DICTrank one, which is the ordering the
base-rate correction follows. The prevalence-matched controls do not follow it ---
they close a similar fraction of the ECE excess in both datasets, 62\% and 66\% ---
so the two analyses are answering different questions and should not be read as
corroborating each other.

Recalibration left discrimination where it was, as monotone transformations must:
AUROC was unchanged to four decimal places under the prior correction. The 0.08 to
0.09 of lost AUROC is a separate problem, and nothing here addresses it.

One bookkeeping point. The uncorrected condition in this analysis gives AUROC
0.660 and ECE 0.447 for DILIrank, against 0.648 and 0.438 in the primary split.
The recalibration pipeline sorts compounds by approval year before fitting, which
changes which rows XGBoost draws under row subsampling. Both values fall inside
the ten-seed range reported below, and every recovery fraction in
Table~\ref{tab:recal} is computed against the uncorrected condition of this same
pipeline.

\subsection{Conformal coverage}

Split conformal prediction offers a distribution-free coverage guarantee that
rests on exchangeability, which a temporal split violates by construction. The
guarantee fails; the interesting part is how (Table~\ref{tab:conformal}).

Marginal coverage fell in both datasets and further in DILIrank. At a nominal
90\%, the DILIrank random reference delivered 90.7\% and the temporal split
73.7\%; at 80\%, 81.0\% against 48.1\%. DICTrank lost less: 89.9\% against 82.4\%
at 90\% nominal, and 80.7\% against 69.7\% at 80\%. That ordering matches the
ordering of the ECE ratios reported above, and is obtained without any
distributional assumption.

Splitting coverage by class shows where the loss falls, and here DICTrank is the
more extreme of the two. Under the DILIrank temporal split the positive class was
covered at 95.8\%, above nominal, while the negative class reached 63.9\% --- a
separation of 32 percentage points, widening to 63 points at nominal 80\% (91.7\%
against 28.7\%). DICTrank separated further still: 99.4\% against 55.3\% at 90\%
nominal, 44 points, and 98.8\% against 23.3\% at 80\%, 75 points. So DICTrank is
the milder case by marginal coverage and the more severe one by class balance,
and reporting only the marginal figure would reverse the impression.

The shortfall falls almost entirely on the negative class in both datasets. In
this setting the negative label is the reassuring one, so a prediction set that
omits it is a model refusing to certify as safe a compound that the label calls
safe. Marginal coverage alone reports a guarantee that is somewhat violated and
hides which half of the guarantee has gone.

Empty prediction sets occurred in one condition only --- 0.3\% of the DICTrank
random reference at nominal 80\% --- and in none of the temporal conditions, so
the shortfall is not an artefact of the quantile running off the end of the
calibration scores. Split conformal guarantees marginal coverage only, so the
class-conditional figures are descriptive; this is why the random reference itself
falls below nominal within the negative class, most visibly at 48.4\% for DICTrank
at nominal 80\%.

\subsection{Label maturity}

Post-marketing labels accumulate, and held-out compounds are younger, so some of
the negative labels in the held-out set may be negative only for now. A gradient
is present in the data over the first 25 years of observation. Among DILIrank
compounds with no more than seven years of observation, 14.0\% carry a positive
label; this rises to 28.6\% at 8--11 years, 52.0\% at 12--15 years and 67.7\% at
16--25 years. DICTrank shows a flatter version of the same trend, from 59.3\% to
78.4\% over the same range.

The gradient does not continue indefinitely. In the oldest stratum, compounds with
more than 25 years of observation --- which is 516 of 796 DILIrank compounds and
472 of 919 DICTrank compounds, the majority of each dataset --- the positive rate
falls back to 64.9\% and 76.1\%. Label accumulation therefore appears to saturate,
and the trend should not be extrapolated as though every negative label will
eventually turn positive. The DICTrank strata sum to 916 rather than 919: three
compounds --- nilotinib, amantadine and selumetinib --- carry an approval year
later than the 2023 label release, which the audit flags as a probable
approval-year matching error rather than a genuine post-release approval. They
have no defined observation window and are therefore absent from this analysis
only. Two of the three, nilotinib and amantadine, are among the corrections tested
in the approval-year sensitivity analysis below; selumetinib is not, because its
earliest matching record is also after 2010 and correcting it would not move it
across the boundary.

Restricting the held-out set to mature labels narrows the DILIrank gap without
closing it. Against a baseline gap of 0.094, requiring at least eight years of
observation gave 0.084, at least ten years 0.077, and at least twelve years 0.072
--- a reduction of 24\% on 50 held-out compounds. Down-weighting negative labels
in proportion to observation time changed the gap by less than 0.003 at every
value of $\tau$. Applying the restriction to training and held-out sets together
widened the gap, to 0.106.

DICTrank did not move in the expected direction at all. Against a baseline of
0.088, the eight-year restriction widened the gap to 0.117 and the ten-year
restriction to 0.093. The largest reduction available under any control was
0.0020, from down-weighting at $\tau = 5$, and the joint restriction widened the
gap to 0.111.

Label maturity therefore accounts for at most about a quarter of the
discrimination gap in DILIrank, on the smallest and most favourable subset
examined, and for essentially none of it in DICTrank. In both datasets the
comparison that holds training and held-out maturity to the same standard widens
rather than narrows the gap. These four controls were evaluated on AUROC only;
calibration was not recomputed under them, so they bear on the discrimination gap
and not on the calibration collapse.

\subsection{Robustness}

Across all 16 combinations of four learners and four representations
(Figure~\ref{fig:robust}a), ECE was worse under the temporal split in
\textbf{16 of 16} cells for both datasets. The median ratio was 2.64 for DILIrank
(range 1.64--5.28) and 1.91 for DICTrank (range 1.43--3.65), and the median
calibration slope 0.184 and 0.172. Discrimination was less uniform. $\Delta$AUROC
was positive in 15 of 16 DILIrank cells, median $+0.096$, and in 16 of 16 DICTrank
cells, median $+0.088$. The one exception, RDKit descriptors with logistic
regression, is also among the weakest cells in the grid under random splitting, so
there was little discrimination there to lose.

Moving the cut-off to 2005 or 2015 preserves the calibration pattern and not the
discrimination pattern (Figure~\ref{fig:robust}b). ECE ratios stayed between 2.71
and 3.67 for DILIrank and between 1.64 and 2.19 for DICTrank, never falling below
1.6 in either dataset. $\Delta$AUROC was the less stable quantity, ranging from
0.041 to 0.097 for DILIrank --- a factor of 2.4 --- and from 0.050 to 0.082 for
DICTrank. Both datasets peak at 2010, but the 2005 and 2015 values swap rank
between them, so the size of the AUROC gap is not a stable function of how far
into the past the boundary is placed. The reliability
share of the Brier degradation was 88.2\%, 86.1\% and 81.4\% for DILIrank at 2005,
2010 and 2015, and 72.9\%, 82.5\% and 90.1\% for DICTrank. The share exceeds 70\%
at every cut-off in both datasets, but it is not constant --- it spans about 7
percentage points in DILIrank and 17 in DICTrank --- and the two datasets move in
opposite directions as the cut-off is made later.

Ten seeds moved AUROC by a standard deviation of 0.0094 for DILIrank and 0.0103
for DICTrank, and ECE by 0.0071 and 0.0074; the seed-42 values reported throughout
sit inside both ranges. The DILIrank AUROC of 0.648 is close to the ten-seed mean
of 0.647; the DICTrank value of 0.620 sits in the upper part of its range
(0.591--0.626, mean 0.611), which is the conservative direction for the claim
being made. Recomputing ECE and the decomposition over 5, 10, 15 and 20 bins under
equal-width and equal-frequency binning moved ECE between 0.424 and 0.454 for
DILIrank and between 0.231 and 0.282 for DICTrank. Reliability exceeded resolution
under every scheme, by a margin of at least 0.19 for DILIrank and 0.04 for
DICTrank.

\subsection{Sensitivity to data-curation decisions}

Three curation decisions could each be argued to drive the result, and each was
tested directly.

\textbf{Approval-year assignment.} An audit against Drugs@FDA identified seven
compounds in the analysis sets assigned to the held-out side of the boundary whose
earliest matching approval record predates 2010: one in DILIrank and six in
DICTrank. Applying those corrections (S1), and additionally excluding phenobarbital
sodium, in clinical use since 1912 but first recorded as an NDA in
2022 (S2), moved the DILIrank ECE ratio from 3.37 to between 3.37 and 3.68 and the
reliability share from 86.1\% to between 84.7\% and 86.2\%. The DICTrank ratio
moved from 1.96 to 2.05 and the share from 82.5\% to 80.3\%. The largest movement
in either dataset is 0.31 in the ratio and 2.2 percentage points in the share.

\textbf{Diagnostic and contrast agents.} Compounds in ATC class V carry safety
labels of a different character from therapeutic agents. Removing them (44
compounds from each dataset, leaving 752 and 875) left the DILIrank ECE ratio at
3.44 against 3.37 and the DICTrank ratio at 1.85 against 1.96, and raised the
reliability share in both datasets, to 91.2\% and 86.7\%. The calibration collapse
is not an artefact of this class.

\textbf{Structure-based deduplication.} Collapsing salts and stereoisomers to a
common parent structure merges entries that occasionally carry opposite labels ---
for instance vancomycin and vancomycin hydrochloride in DICTrank. Rather than
retaining an arbitrary member of such a group, every conflicted group was removed
in full: one group in DILIrank and thirteen in DICTrank, which removes one entry
from the DILIrank analysis set and thirteen from the DICTrank one.
The rule was fixed before its effect was examined, and its effect was not
one-directional: re-running the decomposition on the retained pre-filter sets
(932 and 797 compounds) gives a reliability share of 85.6\% for DILIrank against
86.1\% after the rule, and 84.5\% for DICTrank against 82.5\% --- a rise in one
dataset and a fall in the other. A filter chosen to favour the reported result
would not behave this way.
The pre-filter analysis sets, the re-run output and the full list of removed
groups are given in Additional file~1.

\subsection{Exploratory subgroup analyses}

Held-out compounds sit far from the training set in chemical space. The median
nearest-neighbour Tanimoto similarity to training compounds was 0.271 for
DILIrank, and 115 of 156 held-out compounds (73.7\%) fell below 0.4. Performance
rose with similarity: AUROC 0.543 below 0.3 ($n = 92$), 0.711 between 0.3 and 0.4
($n = 23$), and 0.857 between 0.4 and 0.6 ($n = 29$). The formal tests do not
support the trend. Splitting the held-out set at the median gave a difference of
$+0.068$ with a two-sample bootstrap $p$ of 0.527, and the correlation between
similarity and the predicted probability of the true class was weak
(Spearman $\rho = +0.116$, $p = 0.148$).

The largest subgroup effect was therapeutic rather than structural. Antineoplastic
and immunomodulating agents (ATC class L) accounted for 32 held-out DILIrank
compounds, 84\% of them positive, with an AUROC of 0.289 against 0.649 for the
remaining 118, a difference of $-0.360$ at a bootstrap $p$ of 0.021
(Figure~\ref{fig:atc}). The same comparison in DICTrank gave $-0.048$,
$p = 0.552$, on a subgroup three times larger. Of the eight physicochemical
strata, two reached nominal significance --- DICTrank molecular weight above 500
($p = 0.020$) and DICTrank six or more nitrogen atoms ($p = 0.050$) --- and the
corresponding DILIrank nitrogen stratum points the opposite way.

After Benjamini--Hochberg correction within the exploratory family,
\textbf{none of the twelve tests remains significant}; the smallest $q$ is 0.126,
carried jointly by the DILIrank ATC class L result and the DICTrank molecular
weight stratum. Below-chance AUROC in a 32-compound subgroup is reported for
completeness and supports no claim made in this paper.


%====================================================================%
% Main-text tables                                                   %
% Generated by src/make_tables.py from results/. Do not edit by hand. %
% To change a number, re-run the analysis and then this script.      %
% Journal rules: caption above (<=15 words), legend below, no colour, %
% no thousands separators.                                           %
%====================================================================%

\begin{table}[h]
\caption{Discrimination and calibration under three validation designs}\label{tab:split}
\begin{tabular}{@{}llcc@{}}
\toprule
Dataset & Validation design & AUROC & ECE \\
\midrule
DILIrank & Stratified five-fold CV & 0.7284 $\pm$ 0.0223 & 0.1358 $\pm$ 0.0124 \\
 & Bemis--Murcko scaffold five-fold CV & 0.7306 $\pm$ 0.0183 & 0.1299 $\pm$ 0.0289 \\
 & Temporal split at 2010 & 0.6476 [0.555--0.736] & 0.4384 \\
\midrule
DICTrank & Stratified five-fold CV & 0.6992 $\pm$ 0.0234 & 0.1292 $\pm$ 0.0220 \\
 & Bemis--Murcko scaffold five-fold CV & 0.6775 $\pm$ 0.0355 & 0.1362 $\pm$ 0.0297 \\
 & Temporal split at 2010 & 0.6198 [0.553--0.688] & 0.2702 \\
\botrule
\end{tabular}
\footnotetext{All three designs use the same compounds, model and features within each dataset (796 DILIrank, 919 DICTrank). Cross-validation entries are the mean and standard deviation across five folds; the temporal entry gives the 95\% bootstrap interval for AUROC and a single value for ECE on the 156 and 267 held-out compounds. Scaffold cross-validation changes AUROC by +0.0022 in DILIrank and -0.0217 in DICTrank, that is in opposite directions, and leaves ECE within one standard deviation of stratified cross-validation in both.}
\end{table}

\begin{table}[h]
\caption{Murphy decomposition of the Brier score under random and temporal validation}\label{tab:brier}
\begin{tabular}{@{}lcccc@{}}
\toprule
 & \multicolumn{2}{c}{DILIrank} & \multicolumn{2}{c}{DICTrank} \\
\cmidrule(lr){2-3}\cmidrule(lr){4-5}
Term & Random & Temporal & Random & Temporal \\
\midrule
Uncertainty & 0.2405 & 0.2130 & 0.1982 & 0.2370 \\
Resolution (higher is better) & 0.0528 & 0.0191 & 0.0240 & 0.0125 \\
Reliability (lower is better) & 0.0242 & 0.2327 & 0.0274 & 0.0818 \\
Brier score & 0.2126 & 0.4266 & 0.2046 & 0.3062 \\
\midrule
Reliability share of degradation & \multicolumn{2}{c}{86.1\% [84.6--97.3]} & \multicolumn{2}{c}{82.5\% [75.5--124.0]} \\
Resolution share of degradation & \multicolumn{2}{c}{13.9\%} & \multicolumn{2}{c}{17.5\%} \\
Reliability-to-resolution ratio & \multicolumn{2}{c}{6.19} & \multicolumn{2}{c}{4.72} \\
\midrule
\multicolumn{5}{@{}l}{\textit{Sensitivity: random partitions pooled before binning}} \\
Reliability share of degradation & \multicolumn{2}{c}{90.4\%} & \multicolumn{2}{c}{91.8\%} \\
\botrule
\end{tabular}
\footnotetext{The random reference decomposes each of 20 same-sized partitions separately and averages the terms, so both columns rest on the same held-out size (156 and 267). The standard deviation across those partitions is 0.0066 and 0.0074 for reliability and 0.0132 and 0.0086 for resolution. \textbf{The uncertainty row comes from a separate 200-partition run} (the R1 control of Table~\ref{tab:prev}), because the 20-partition output does not retain it; the random column of the identity therefore closes only to 0.2119 and 0.2016 rather than exactly. Shares are the change in each model-dependent term divided by the sum of both changes; intervals are 2000-fold bootstrap percentiles. An upper bound above 100\% arises in resamples where resolution improves rather than deteriorates, which occurred in 0.9\% of DILIrank and 26.9\% of DICTrank resamples. Uncertainty is fixed by the held-out class balance and is not attributed. In the temporal columns, where an exact Brier score on the same predictions is available, the binned reconstruction and the exact score agree to within 0.002 (0.4271 against 0.4266 and 0.3050 against 0.3062). The pooled row gives the random reference twenty times the sample of the temporal column and is reported only to show the direction of that bias.}
\end{table}

\begin{table}[h]
\caption{Prevalence-matched random controls for expected calibration error}\label{tab:prev}
\begin{tabular}{@{}llccc@{}}
\toprule
Dataset & Control & ECE & Ratio & $p$ \\
\midrule
DILIrank & Temporal split at 2010 & 0.4384 [0.374--0.527] & --- & --- \\
 & (R1) Unconstrained random & 0.1309 $\pm$ 0.0285 & 3.35 & 0.0050 \\
 & (R2) Held-out rate matched & 0.2495 $\pm$ 0.0322 & 1.76 & 0.0050 \\
 & (R3) Both rates matched & 0.3211 $\pm$ 0.0316 & 1.37 & 0.0050 \\
\midrule
DICTrank & Temporal split at 2010 & 0.2702 [0.228--0.336] & --- & --- \\
 & (R1) Unconstrained random & 0.1299 $\pm$ 0.0237 & 2.08 & 0.0050 \\
 & (R2) Held-out rate matched & 0.1866 $\pm$ 0.0250 & 1.45 & 0.0050 \\
 & (R3) Both rates matched & 0.2226 $\pm$ 0.0233 & 1.21 & 0.0348 \\
\botrule
\end{tabular}
\footnotetext{Each control is averaged over 200 partitions. R2 matches the held-out positive rate to the temporal held-out set; R3 matches the training rate as well, at the cost of a smaller held-out set (92 and 189 compounds against 156 and 267). Ratio is temporal ECE divided by control ECE. One-sided $p$ uses the bias-corrected estimator $(1+r)/(1+N)$, where $r$ is the number of partitions with ECE at or above the temporal value; $r$ was 0 in 5 of the six comparisons and 6 (DICTrank R3) in the remaining one. The value 0.0050 is the lower bound available at 200 repeats, so 5 of these $p$ values are at that bound and the true values are smaller. All six survive Benjamini--Hochberg correction within this family, the largest $q$ being 0.0348.}
\end{table}

\begin{table}[h]
\caption{Recalibration strategies evaluated on the temporal held-out set}\label{tab:recal}
\begin{tabular}{@{}lcccc@{}}
\toprule
 & \multicolumn{2}{c}{DILIrank} & \multicolumn{2}{c}{DICTrank} \\
\cmidrule(lr){2-3}\cmidrule(lr){4-5}
Strategy & ECE & Gap closed & ECE & Gap closed \\
\midrule
(A) No recalibration & 0.4466 & --- & 0.2726 & --- \\
(B) Platt scaling, internal & 0.4147 & 8\% & 0.1679 & 42\% \\
(C) Isotonic regression, internal & 0.4085 & 10\% & 0.1739 & 40\% \\
(D) Prior-probability correction & 0.2959 & 38\% & 0.2438 & 12\% \\
(E) Oracle Platt scaling & 0.0482 & 100\% & 0.0234 & 100\% \\
\midrule
(B)--(C) across calibration fractions & \multicolumn{2}{c}{1.3\%--20.0\%} & \multicolumn{2}{c}{35.0\%--54.4\%} \\
\botrule
\end{tabular}
\footnotetext{Strategies (B) and (C) fit on the most recent 25\% of training compounds by approval year (160 and 163 compounds, spanning 1997--2009 and 1998--2009) and use no held-out information. Strategy (D) rescales the intercept of the full-training model using the held-out positive rate; it is not prospective and is reported to isolate the pure base-rate component. Table~\ref{tab:prior} shows that estimating that rate without held-out labels does not work, so (D) belongs with (E) as a ceiling. Strategy (E) is fitted on the held-out labels. Gap closed is the fraction of the distance between (A) and (E) that a strategy recovers. \textbf{The last row is the range over calibration fractions 0.15--0.40 and is the quantity to quote}; the single-fraction values above are one point in that range. (D) and (E) are monotone transformations of (A) and leave its AUROC unchanged at 0.6603 and 0.6045; (B) and (C) are applied to a model refitted without the calibration set, so their AUROC differs (0.6559 and 0.6403; 0.6131 and 0.6153), and isotonic regression can move AUROC by as much as 0.0613 at the smallest calibration fraction. Row (A) differs from the primary split because this pipeline sorts by approval year before fitting, which changes the rows drawn under subsampling.}
\end{table}

\begin{table}[h]
\caption{Split conformal coverage under random and temporal validation}\label{tab:conformal}
\begin{tabular}{@{}llcccc@{}}
\toprule
 &  & \multicolumn{2}{c}{DILIrank} & \multicolumn{2}{c}{DICTrank} \\
\cmidrule(lr){3-4}\cmidrule(lr){5-6}
Nominal & Coverage & Random & Temporal & Random & Temporal \\
\midrule
90\% & Marginal & 0.907 & 0.737 & 0.899 & 0.824 \\
 & Positive class & 0.953 & 0.958 & 0.982 & 0.994 \\
 & Negative class & 0.842 & 0.639 & 0.696 & 0.553 \\
\midrule
80\% & Marginal & 0.810 & 0.481 & 0.807 & 0.697 \\
 & Positive class & 0.890 & 0.917 & 0.939 & 0.988 \\
 & Negative class & 0.697 & 0.287 & 0.484 & 0.233 \\
\botrule
\end{tabular}
\footnotetext{The random reference is averaged over partitions with identical fitting, calibration and held-out sizes (calibration sets of 160 and 163, held-out sets of 156 and 267). Empty prediction sets occurred in 1 of the eight conditions, all on the random side: DICTrank at nominal 80\% at a rate of 0.003. Every temporal condition gave none. Split conformal guarantees marginal coverage only; the class-conditional rows are descriptive, which is why the random reference can itself fall below nominal within a class. The marginal shortfall under the temporal split is 0.163 and 0.076 at nominal 90\% and 0.319 and 0.103 at nominal 80\%, and in every case it falls on the negative class.}
\end{table}

\begin{table}[h]
\caption{Reversed temporal split as a control against label censoring}\label{tab:reverse}
\begin{tabular}{@{}llcccccccc@{}}
\toprule
Dataset & Direction and control & Train $n$ & Test $n$ & ECE & Slope & Ratio & Reliability share & Valid reps & $p$ \\
\midrule
DILIrank & Forward, temporal split & 640 & 156 & 0.4384 & 0.2700 & --- & --- & --- & --- \\
 & Reverse, temporal split & 156 & 640 & 0.4857 & 0.1331 & --- & --- & --- & --- \\
 & Reverse, (R1) unconstrained random & 156 & 640 & 0.1990 & 0.2406 & 2.44 & 95.0\% & 200/200 & 0.0050 \\
 & Reverse, (R2) held-out rate matched & 156 & 573 & 0.1911 & 0.2423 & 2.54 & 95.5\% & 200/200 & 0.0050 \\
 & Reverse, (R3) both rates matched & 156 & 583 & 0.3916 & 0.1745 & 1.24 & 96.0\% & 40/200 & 0.0244 \\
\midrule
DICTrank & Forward, temporal split & 652 & 267 & 0.2702 & 0.2223 & --- & --- & --- & --- \\
 & Reverse, temporal split & 267 & 652 & 0.2060 & 0.2384 & --- & --- & --- & --- \\
 & Reverse, (R1) unconstrained random & 267 & 652 & 0.1589 & 0.2714 & 1.30 & 86.1\% & 200/200 & 0.0100 \\
 & Reverse, (R2) held-out rate matched & 267 & 614 & 0.1443 & 0.2686 & 1.43 & 91.1\% & 200/200 & 0.0050 \\
 & Reverse, (R3) both rates matched & 267 & 614 & 0.1821 & 0.2517 & 1.13 & 76.0\% & 200/200 & 0.1294 \\
\botrule
\end{tabular}
\footnotetext{Training on compounds approved from 2010 onwards and testing on earlier ones makes the held-out labels the mature ones, so label immaturity cannot explain a collapse in this direction. Calibration still degrades: against a size-matched random control the ratio is 2.44 and 1.30 ($p$ = 0.0050 and 0.0100), and the reliability share, 95.0\% and 86.1\%, is higher than in the forward direction. \textbf{Three caveats belong with this table.} Under the strictest control (R3) the two datasets part company, DILIrank at 1.24 ($p$ = 0.0244) against DICTrank at 1.13 ($p$ = 0.1294, not significant), so the reversed evidence is weaker than the forward evidence of Table~\ref{tab:prev}. The DILIrank R3 control retains only 40 of 200 repeats, because matching its training positive rate of 0.3077 forces a subsample that usually fails the minimum-size guard; its $p$ therefore sits at the 0.0244 bound of 40 valid repeats rather than the 0.0050 bound used elsewhere, and the repeat count must be quoted with it. Absolute AUROC is not comparable with the forward direction because the reversed training set holds only 156 and 267 compounds, whereas the ECE ratio is comparable because each control uses the same training size. This is a control experiment, not a deployment scenario.}
\end{table}

\begin{table}[h]
\caption{Non-temporal splits that saturate covariate shift}\label{tab:nontemporal}
\begin{tabular}{@{}llcccccccc@{}}
\toprule
Dataset & Split & Test $n$ & Prevalence shift & Membership AUC & AUROC & ECE & Control ECE & Ratio & $p$ \\
\midrule
DILIrank & Temporal at 2010 & 156 & -0.3470 & 0.6835 & 0.6476 & 0.4384 & 0.2495 & 1.76 & 0.0050 \\
 & Molecular weight, upper half & 156 & -0.0679 & 1.0000 & 0.6889 & 0.1868 & 0.1431 & 1.30 & 0.0746 \\
 & Calculated logP, upper half & 156 & +0.0596 & 1.0000 & 0.6592 & 0.1789 & 0.1282 & 1.40 & 0.0348 \\
 & Fingerprint cluster hold-out & 156 & +0.0198 & 0.9827 & 0.7195 & 0.1570 & 0.1284 & 1.22 & 0.1741 \\
 & Random & 156 & -0.0600 & 0.4761 & 0.7068 & 0.1340 & 0.1410 & 0.95 & 0.5672 \\
\midrule
DICTrank & Temporal at 2010 & 267 & -0.1557 & 0.7917 & 0.6198 & 0.2702 & 0.1866 & 1.45 & 0.0050 \\
 & Molecular weight, upper half & 267 & -0.0712 & 0.9993 & 0.5528 & 0.2011 & 0.1529 & 1.31 & 0.0249 \\
 & Calculated logP, upper half & 267 & +0.0079 & 0.9990 & 0.5901 & 0.1726 & 0.1284 & 1.34 & 0.0299 \\
 & Fingerprint cluster hold-out & 267 & +0.0132 & 0.9880 & 0.7081 & 0.1255 & 0.1262 & 0.99 & 0.5075 \\
 & Random & 267 & +0.0449 & 0.4532 & 0.6834 & 0.1237 & 0.1190 & 1.04 & 0.4030 \\
\botrule
\end{tabular}
\footnotetext{Each split keeps the held-out size of the temporal split and carries its own random control matched to that split's held-out prevalence (200 repeats, the R2 design of Table~\ref{tab:prev}). Membership AUC is a five-fold cross-validated classifier for which half a compound belongs to, computed by the same procedure as the era classifier described in Methods. \textbf{The design did not achieve matched shift and is reported as an upper-bound contrast instead.} Molecular weight and logP are model inputs, so splitting on them drives membership AUC to 1.0000, and the fingerprint-cluster split reaches 0.9827 and 0.9880, against 0.6835 and 0.7917 for the temporal split. Even with shift saturated, no non-temporal ratio reaches the temporal one: the six shift-saturated ratios span 0.99--1.40, and including the two random splits the eight non-temporal ratios span 0.95--1.40, while the temporal split gives 1.76 and 1.45. The random split, with membership AUC near chance, gives 0.95 and 1.04, which confirms the measurement behaves. Across the ten rows the ratio tracks the absolute prevalence shift (Spearman rho = +0.515) more closely than membership AUC (rho = +0.255); \textbf{with ten rows no significance is claimed for either correlation and they are to be read as an ordering only}. The row T reproduces the R2 ratio of Table~\ref{tab:prev} exactly, which cross-checks the new script against the existing pipeline.}
\end{table}

\begin{table}[h]
\caption{Estimating the held-out prior without held-out labels}\label{tab:prior}
\begin{tabular}{@{}lllcccc@{}}
\toprule
Dataset & Prior estimator & Information used & $\hat{\pi}$ & Error & ECE & Gap closed \\
\midrule
DILIrank & (P0) Assume no shift: training rate & training labels only & 0.6547 & +0.3470 & 0.4466 & 0.0\% \\
 & (P1) Mean predicted probability & none & 0.7411 & +0.4334 & 0.4832 & -9.2\% \\
 & (P2) Expectation--maximisation & none & 0.8235 & +0.5158 & 0.5158 & -17.4\% \\
 & (P3) Black-box shift estimation & training labels only & 0.5645 & +0.2568 & 0.4244 & 5.6\% \\
 & (P4) Oracle: held-out rate & held-out labels & 0.3077 & +0.0000 & 0.2959 & 37.8\% \\
\midrule
DICTrank & (P0) Assume no shift: training rate & training labels only & 0.7699 & +0.1557 & 0.2726 & 0.0\% \\
 & (P1) Mean predicted probability & none & 0.8443 & +0.2301 & 0.2997 & -10.9\% \\
 & (P2) Expectation--maximisation & none & 0.9550 & +0.3408 & 0.3490 & -30.7\% \\
 & (P3) Black-box shift estimation & training labels only & 1.0000 & +0.3858 & 0.3858 & -45.4\% \\
 & (P4) Oracle: held-out rate & held-out labels & 0.6142 & +0.0000 & 0.2438 & 11.5\% \\
\botrule
\end{tabular}
\footnotetext{Only the prior estimate changes between rows; the correction applied to the intercept is the one of strategy (D) in Table~\ref{tab:recal}, and (P4) is that strategy. Gap closed is measured against the same uncorrected and oracle endpoints as Table~\ref{tab:recal}. \textbf{The result is a clear negative.} Every estimator that does not see held-out labels errs upward, by +0.1557 to +0.5158, and 5 of the 8 estimates make calibration worse rather than better. The best label-free result is 5.6\% for DILIrank and 0.0\% for DICTrank, against 37.8\% and 11.5\% for the oracle; for DICTrank the best available action is to do nothing. The common direction of the error is the point: all four estimators start from the model's own scores or hard predictions, and the temporal split has already pushed those upward (mean predicted probability 0.7321 and 0.8453), so a score-based estimate reads more positives while the true prevalence has fallen. Estimating a prior from a model whose calibration has collapsed is circular. In DICTrank the black-box confusion matrix becomes near-singular (determinant 0.0050) and the estimate saturates at 1.0000. Deployable recalibration is therefore limited to strategies (B) and (C) of Table~\ref{tab:recal}.}
\end{table}

\begin{table}[h]
\caption{Concept shift isolated by holding training size and prevalence fixed}\label{tab:concept}
\begin{tabular}{@{}llccccccccc@{}}
\toprule
Dataset & Arm & Train era & Train $n$ & Train rate & AUROC & ECE & Slope & Reliability & Ratio to A & $p$ \\
\midrule
DILIrank & (A) Era-matched, random 5-fold CV & post-2010 & 125 & 0.3077 & 0.7290 & 0.1443 & 0.4205 & 0.0305 & --- & --- \\
 & (A2) Era-matched, scaffold 5-fold CV & post-2010 & 125 & 0.3077 & 0.6925 & 0.1553 & 0.3533 & 0.0276 & --- & --- \\
 & (B) Forward, full pre-2010 training set & pre-2010 & 640 & 0.6547 & 0.6476 & 0.4384 & 0.2700 & 0.2327 & --- & --- \\
 & (C) Forward, training size matched to A & pre-2010 & 125 & 0.6550 & 0.6148 & 0.4643 & 0.1918 & 0.2526 & 3.22 & 0.0196 \\
 & (D) Forward, size and prevalence matched to A & pre-2010 & 125 & 0.3040 & 0.6177 & 0.2231 & 0.1953 & 0.0751 & 1.55 & 0.0392 \\
\midrule
DICTrank & (A) Era-matched, random 5-fold CV & post-2010 & 214 & 0.6142 & 0.5822 & 0.2165 & 0.1535 & 0.0629 & --- & --- \\
 & (A2) Era-matched, scaffold 5-fold CV & post-2010 & 214 & 0.6142 & 0.5636 & 0.2496 & 0.1241 & 0.0743 & --- & --- \\
 & (B) Forward, full pre-2010 training set & pre-2010 & 652 & 0.7699 & 0.6198 & 0.2702 & 0.2223 & 0.0818 & --- & --- \\
 & (C) Forward, training size matched to A & pre-2010 & 214 & 0.7746 & 0.5705 & 0.2908 & 0.1337 & 0.0939 & 1.34 & 0.0196 \\
 & (D) Forward, size and prevalence matched to A & pre-2010 & 214 & 0.6121 & 0.5771 & 0.2380 & 0.1467 & 0.0715 & 1.10 & 0.2549 \\
\botrule
\end{tabular}
\footnotetext{Every arm is scored on the same held-out compounds, the 156 and 267 approved from 2010 onwards; arms A and A2 use out-of-fold predictions so no compound is scored by a model trained on it. Arms A and D differ only in the era the training compounds come from, their size and positive rate being matched, so the gap between them is the concept-shift term. It is +0.0788 for DILIrank ($p$ = 0.0392, 1 of 50 controls matched arm A) and +0.0216 for DICTrank ($p$ = 0.2549, 12 of 50), so \textbf{concept shift is demonstrated for hepatotoxicity and not for cardiotoxicity}. Arm A2 separates Bemis--Murcko scaffolds across folds and costs only +0.0110 and +0.0331 of ECE, which rules out analogue leakage as the explanation for arm A. Arm B reproduces the temporal split of Table~\ref{tab:split} to four decimals, which fixes this experiment on the same pipeline. Along the path A to D to C to B the three increments sum exactly to +0.2941 and +0.0537: era +0.0788 and +0.0216, training prevalence +0.2412 and +0.0527, training size -0.0258 and -0.0206. \textbf{The increments depend on the path order}, which distributes the interactions, and the negative size term means that enlarging the training set from the matched size to the full pre-2010 set slightly improved calibration, so small training sets do not explain the collapse. The $p$ values compare each control against a reference distribution of 50 resamples and are not exchangeable permutation tests; read the ECE differences first.}
\end{table}

\begin{table}[h]
\caption{Calibration tracks the training prevalence gap while discrimination does not}\label{tab:sweep}
\begin{tabular}{@{}lcccccccc@{}}
\toprule
Dataset & Training rate & Gap & Train $n$ & AUROC & ECE & ECE SD & Slope & Reliability \\
\midrule
DILIrank & 0.2000 & -0.1077 & 262 & 0.6392 & 0.2153 & 0.0205 & 0.2656 & 0.0544 \\
 & 0.3000 & -0.0077 & 262 & 0.6629 & 0.1752 & 0.0284 & 0.3036 & 0.0419 \\
 & 0.3077$^{a}$ & +0.0000 & 262 & 0.6507 & 0.1827 & 0.0308 & 0.2750 & 0.0464 \\
 & 0.4000 & +0.0923 & 262 & 0.6668 & 0.1813 & 0.0349 & 0.3025 & 0.0531 \\
 & 0.5000 & +0.1923 & 262 & 0.6736 & 0.2238 & 0.0328 & 0.3114 & 0.0788 \\
 & 0.6000 & +0.2923 & 262 & 0.6758 & 0.3034 & 0.0288 & 0.3153 & 0.1288 \\
 & 0.6547$^{b}$ & +0.3470 & 262 & 0.6635 & 0.3602 & 0.0291 & 0.2910 & 0.1683 \\
 & 0.7000 & +0.3923 & 262 & 0.6589 & 0.4023 & 0.0267 & 0.2872 & 0.1976 \\
 & 0.8000 & +0.4923 & 262 & 0.6434 & 0.5079 & 0.0257 & 0.2632 & 0.2836 \\
\midrule
DICTrank & 0.2000 & -0.4142 & 178 & 0.6256 & 0.4441 & 0.0246 & 0.2303 & 0.2172 \\
 & 0.3000 & -0.3142 & 178 & 0.6316 & 0.3544 & 0.0278 & 0.2310 & 0.1531 \\
 & 0.4000 & -0.2142 & 178 & 0.6443 & 0.2681 & 0.0322 & 0.2509 & 0.0969 \\
 & 0.5000 & -0.1142 & 178 & 0.6424 & 0.2190 & 0.0298 & 0.2474 & 0.0650 \\
 & 0.6000 & -0.0142 & 178 & 0.6398 & 0.2045 & 0.0277 & 0.2370 & 0.0537 \\
 & 0.6142$^{a}$ & +0.0000 & 178 & 0.6398 & 0.2001 & 0.0310 & 0.2429 & 0.0513 \\
 & 0.7000 & +0.0858 & 178 & 0.6387 & 0.2189 & 0.0232 & 0.2436 & 0.0564 \\
 & 0.7699$^{b}$ & +0.1557 & 178 & 0.6245 & 0.2517 & 0.0238 & 0.2257 & 0.0717 \\
 & 0.8000 & +0.1858 & 178 & 0.6222 & 0.2697 & 0.0192 & 0.2215 & 0.0808 \\
\botrule
\end{tabular}
\footnotetext{Random splits only; no temporal information enters this experiment. The held-out set is fixed at the size and positive rate of the temporal test set, marked $^{a}$, and the training set is held at a fixed size while only its positive rate is varied; $^{b}$ marks the pre-2010 rate, so that row carries the gap the temporal split actually imposes. Each row averages 100 repeats. \textbf{ECE varies by a factor of 2.90 and 2.22 across the sweep while AUROC stays within 0.0366 and 0.0221}, which is the paper's central claim reproduced in a single controlled experiment. ECE is minimised at a gap of -0.0077 and +0.0000, that is at zero within one grid step, in both datasets. Regressed on the absolute gap, ECE rises by 0.0646 [0.063--0.066] and 0.0576 [0.056--0.059] per 0.10 of gap. The two sides are not symmetric: for DILIrank the slopes are 0.0401 below and 0.0681 above the held-out rate, against 0.0571 and 0.0501 for DICTrank, so overstating the training prevalence is the more damaging error for hepatotoxicity only. The calibration slope stays near 0.29 and 0.24 throughout while reliability rises 6.77-fold and 4.23-fold, which locates the damage in the intercept rather than the spread of the predictions. \textbf{This sweep should be read for shape and slope, not for shares}: its training set is held at 262 and 178 compounds, far below the 640 and 652 of the temporal split, which raises the zero-gap baseline and compresses the range. The shares belong to Table~\ref{tab:prev} and Table~\ref{tab:concept}.}
\end{table}

\begin{table}[h]
\caption{Recency weighting of the training set does not repair calibration}\label{tab:recency}
\begin{tabular}{@{}llcccccccccc@{}}
\toprule
Dataset & Training scheme & Train $n$ & Effective $n$ & Train rate & Gap & AUROC & ECE & 95\% CI & Slope & Reliability & Resolution \\
\midrule
DILIrank & Uniform weights (= temporal split) & 640 & 640.0 & 0.6547 & +0.3470 & 0.6476 & 0.4384 & [0.374--0.527] & 0.2700 & 0.2327 & 0.0191 \\
 & Half-life 30 years & 640 & 553.7 & 0.6893 & +0.3816 & 0.6501 & 0.4424 & [0.378--0.534] & 0.2597 & 0.2378 & 0.0172 \\
 & Half-life 20 years & 640 & 486.9 & 0.6995 & +0.3918 & 0.6493 & 0.4448 & [0.378--0.535] & 0.2590 & 0.2367 & 0.0094 \\
 & Half-life 15 years & 640 & 426.4 & 0.7057 & +0.3980 & 0.6399 & 0.4567 & [0.384--0.541] & 0.2433 & 0.2415 & 0.0141 \\
 & Half-life 10 years & 640 & 333.3 & 0.7096 & +0.4019 & 0.6458 & 0.4615 & [0.395--0.547] & 0.2674 & 0.2464 & 0.0151 \\
 & Half-life 7 years & 640 & 255.0 & 0.7047 & +0.3970 & 0.6508 & 0.4543 & [0.390--0.542] & 0.2970 & 0.2440 & 0.0202 \\
 & Half-life 5 years & 640 & 189.6 & 0.6914 & +0.3838 & 0.6375 & 0.4587 & [0.387--0.536] & 0.3130 & 0.2364 & 0.0240 \\
 & Half-life 3 years & 640 & 111.9 & 0.6590 & +0.3513 & 0.6586 & 0.4385 & [0.367--0.514] & 0.3998 & 0.2134 & 0.0207 \\
 & Truncated to last 30 years & 367 & 367.0 & 0.7384 & +0.4307 & 0.6173 & 0.5121 & [0.441--0.595] & 0.2214 & 0.2937 & 0.0101 \\
 & Truncated to last 20 years & 272 & 272.0 & 0.7574 & +0.4497 & 0.6512 & 0.5180 & [0.447--0.593] & 0.2832 & 0.2955 & 0.0251 \\
 & Truncated to last 15 years & 221 & 221.0 & 0.7421 & +0.4344 & 0.6698 & 0.5047 & [0.431--0.582] & 0.3058 & 0.2862 & 0.0167 \\
 & Truncated to last 10 years & 124 & 124.0 & 0.6774 & +0.3697 & 0.6466 & 0.4693 & [0.401--0.554] & 0.2748 & 0.2566 & 0.0148 \\
\midrule
DICTrank & Uniform weights (= temporal split) & 652 & 652.0 & 0.7699 & +0.1557 & 0.6198 & 0.2702 & [0.225--0.335] & 0.2223 & 0.0818 & 0.0125 \\
 & Half-life 30 years & 652 & 573.7 & 0.7837 & +0.1695 & 0.5975 & 0.2809 & [0.232--0.343] & 0.1956 & 0.0843 & 0.0066 \\
 & Half-life 20 years & 652 & 512.7 & 0.7884 & +0.1741 & 0.5712 & 0.2664 & [0.225--0.335] & 0.1573 & 0.0830 & 0.0034 \\
 & Half-life 15 years & 652 & 455.9 & 0.7918 & +0.1775 & 0.5557 & 0.2810 & [0.232--0.342] & 0.1244 & 0.0909 & 0.0098 \\
 & Half-life 10 years & 652 & 365.0 & 0.7957 & +0.1815 & 0.5343 & 0.2804 & [0.233--0.343] & 0.0807 & 0.0912 & 0.0038 \\
 & Half-life 7 years & 652 & 284.1 & 0.7979 & +0.1836 & 0.5047 & 0.2818 & [0.232--0.341] & 0.0242 & 0.0925 & 0.0056 \\
 & Half-life 5 years & 652 & 213.9 & 0.7994 & +0.1852 & 0.4792 & 0.2679 & [0.220--0.332] & -0.0395 & 0.0841 & 0.0050 \\
 & Half-life 3 years & 652 & 129.0 & 0.8053 & +0.1911 & 0.5163 & 0.2258 & [0.186--0.294] & 0.0438 & 0.0657 & 0.0026 \\
 & Truncated to last 30 years & 414 & 414.0 & 0.8043 & +0.1901 & 0.5413 & 0.3197 & [0.268--0.379] & 0.0844 & 0.1141 & 0.0084 \\
 & Truncated to last 20 years & 309 & 309.0 & 0.8058 & +0.1916 & 0.4938 & 0.3252 & [0.277--0.390] & -0.0147 & 0.1178 & 0.0069 \\
 & Truncated to last 15 years & 252 & 252.0 & 0.7897 & +0.1755 & 0.4933 & 0.3171 & [0.271--0.383] & -0.0079 & 0.1204 & 0.0101 \\
 & Truncated to last 10 years & 143 & 143.0 & 0.7692 & +0.1550 & 0.4870 & 0.2748 & [0.230--0.337] & -0.0185 & 0.0957 & 0.0034 \\
 & Truncated to last 5 years & 70 & 70.0 & 0.7857 & +0.1715 & 0.5267 & 0.2816 & [0.237--0.352] & 0.0582 & 0.0931 & 0.0096 \\
\botrule
\end{tabular}
\footnotetext{Weights are $0.5^{(2010-\mathrm{year})/\tau}$, so a smaller half-life $\tau$ concentrates the training set on recent approvals; uniform weights reproduce the temporal split of Table~\ref{tab:split}. Truncation keeps only the last $K$ years and is the simpler alternative a practitioner would try first. Effective $n$ is $(\sum w)^2/\sum w^2$ and shows what the weighting costs. \textbf{For DILIrank no half-life beats uniform weighting}: the best ECE over the whole sweep, 0.4384, is the uniform value itself, and truncation is worse at every window. \textbf{For DICTrank the apparent gain is bought by destroying resolution}: at $\tau$ = 3 years ECE falls from 0.2702 to 0.2258, but Murphy resolution falls from 0.0125 to 0.0026, roughly a fifth of its value, AUROC falls from 0.6198 to 0.5163 and the calibration slope from 0.2223 to 0.0438, on an effective sample of 129 rather than 652. Flattening the predictions lowers ECE and leaves a model that no longer separates the classes, which is the very error this paper argues against. Half-lives were chosen without looking at the held-out set, by an inner split of the training era (fit before 2000, select on 2000--2009); that selection rule did find the sweep optimum for DICTrank, so the procedure works even though what it selects is not useful. Read with Table~\ref{tab:recal}: the only intervention that recovers a large share of the lost calibration uses labelled recent compounds.}
\end{table}

\begin{table}[h]
\caption{Calibration degrades across the whole chemical space, not only unfamiliar compounds}\label{tab:applicability}
\begin{tabular}{@{}llcccccccc@{}}
\toprule
Dataset & Stratum & $n$ & Prevalence & AUROC & ECE & Control ECE & Ratio & 95\% CI & $p$ \\
\midrule
DILIrank & Nearest-neighbour Tanimoto, lowest third & 56 & 0.4107 & 0.5731 & 0.4254 & 0.2368 & 1.80 & [1.320--2.377] & 0.0323 \\
 & Nearest-neighbour Tanimoto, middle third & 48 & 0.2292 & 0.4816 & 0.5897 & 0.3429 & 1.72 & [1.393--2.078] & 0.0323 \\
 & Nearest-neighbour Tanimoto, highest third & 52 & 0.2692 & 0.8214 & 0.3764 & 0.3211 & 1.17 & [0.872--1.575] & 0.1935 \\
 & Scaffold seen in training set & 34 & --- & --- & --- & --- & --- & --- & underpowered \\
 & Scaffold unseen in training set & 122 & 0.3361 & 0.6022 & 0.4518 & 0.2622 & 1.72 & [1.480--2.096] & 0.0323 \\
\midrule
DICTrank & Nearest-neighbour Tanimoto, lowest third & 90 & 0.6111 & 0.5512 & 0.3085 & 0.2102 & 1.47 & [1.063--1.952] & 0.0323 \\
 & Nearest-neighbour Tanimoto, middle third & 88 & 0.5909 & 0.6031 & 0.3123 & 0.2173 & 1.44 & [1.048--1.896] & 0.0323 \\
 & Nearest-neighbour Tanimoto, highest third & 89 & 0.6404 & 0.7171 & 0.2724 & 0.1852 & 1.47 & [1.012--2.015] & 0.0323 \\
 & Scaffold seen in training set & 50 & 0.6000 & 0.7883 & 0.2820 & 0.2457 & 1.15 & [0.684--1.661] & 0.1613 \\
 & Scaffold unseen in training set & 217 & 0.6175 & 0.5820 & 0.2813 & 0.2051 & 1.37 & [1.102--1.726] & 0.0323 \\
\botrule
\end{tabular}
\footnotetext{Each stratum is compared against random splits matched to its own size and positive rate, because the prevalence differences documented in Table~\ref{tab:sweep} confound a raw comparison across strata; underpowered strata follow the rule used for the disjoint subsets ($n<40$ or minority class $<10$). \textbf{Of the 9 strata that could be computed, 7 have a ratio whose confidence interval lies entirely above one}, so the time split is worse than a matched random split whether the held-out compounds are chemically familiar or not. The similarity gradient, lowest third minus highest third of the ratio, is 0.63 [-0.037, 1.266] for DILIrank and -0.01 [-0.695, 0.661] for DICTrank, both intervals covering zero, while the discrimination gradient, highest third minus lowest third of AUROC, is 0.2484 [0.024--0.451] and 0.1656 [0.002--0.334], neither covering zero. \textbf{This pair of results must not be read as proof that the two measures behave differently}: contrasting a significant with a non-significant interval is not a test of their difference, the two quantities are on different scales, and the calibration interval still admits a gradient as large as 1.27. What the data support is narrower: discrimination varies with similarity, and no variation in the calibration ratio was detectable at these stratum sizes. The calibration intervals resample only the held-out compounds and hold the control mean fixed, so they are narrower than the true intervals.}
\end{table}

\begin{table}[h]
\caption{Under temporal validation the model is outscored by a constant}\label{tab:trivial}
\begin{tabular}{@{}llcccc@{}}
\toprule
Dataset & Predictor & Brier (temporal) & BSS (temporal) & Brier (random) & BSS (random) \\
\midrule
DILIrank & Trained model (XGBoost) & 0.4271 & -0.2809 & 0.2302 & +0.0868 \\
 & Constant = training prevalence & 0.3334 & +0.0000 & 0.2521 & +0.0000 \\
 & Constant = 0.5 & 0.2500 & +0.2502 & 0.2500 & +0.0084 \\
 & Constant = validation prevalence (oracle) & 0.2130 & +0.3611 & 0.2485 & +0.0143 \\
\midrule
DICTrank & Trained model (XGBoost) & 0.3050 & -0.1676 & 0.1811 & +0.0272 \\
 & Constant = training prevalence & 0.2612 & +0.0000 & 0.1862 & +0.0000 \\
 & Constant = 0.5 & 0.2500 & +0.0429 & 0.2500 & -0.3427 \\
 & Constant = validation prevalence (oracle) & 0.2370 & +0.0928 & 0.1842 & +0.0108 \\
\botrule
\end{tabular}
\footnotetext{Brier score, lower is better. The Brier Skill Score is $1-\mathrm{BS}_{\mathrm{model}}/\mathrm{BS}_{\mathrm{C0}}$, so a negative value means the model is worse than predicting the training prevalence for every compound. C0 is the only one of the three constants that can be computed at deployment time; C2 uses the validation labels and is an unreachable reference. \textbf{Under the temporal split the model is beaten by 6 of the 6 dataset-constant pairs, and under a random split of the same compounds it beats 6 of 6}, with features, model, hyperparameters and seed held fixed. AUROC on the same temporal split is 0.6476 and 0.6198, so ranking survives while the probabilities do not. Constants have no ranking and no AUROC.}
\end{table}

\begin{table}[h]
\caption{Calibration-error ratios exclude one in every control design}\label{tab:ratioci}
\begin{tabular}{@{}llcccc@{}}
\toprule
Dataset & Control design & Control ECE & Ratio & 95\% CI & ECE difference 95\% CI \\
\midrule
DILIrank & Random (standard) & 0.1309 & 3.35 & [2.878--4.005] & [0.246--0.392] \\
 & Random, validation prevalence matched & 0.2495 & 1.76 & [1.511--2.097] & [0.127--0.273] \\
 & Random, both prevalences matched & 0.3211 & 1.37 & [1.174--1.628] & [0.056--0.202] \\
\midrule
DICTrank & Random (standard) & 0.1299 & 2.08 & [1.754--2.594] & [0.098--0.207] \\
 & Random, validation prevalence matched & 0.1866 & 1.45 & [1.225--1.806] & [0.042--0.150] \\
 & Random, both prevalences matched & 0.2226 & 1.21 & [1.024--1.516] & [0.005--0.114] \\
\botrule
\end{tabular}
\footnotetext{The ratio is the temporal-split ECE divided by the mean ECE of 200 matched random-split controls. Intervals are 2.5--97.5 percentiles of 4,000 bootstrap resamples of the validation set; within each resample the control repetitions are resampled as well, so both the numerator and the denominator contribute sampling error. \textbf{All six lower bounds exceed one}, the smallest being 1.024 for DICTrank under the design that matches training and validation prevalence simultaneously. The ECE differences likewise exclude zero. Including the denominator's variability widens the intervals by a factor of at most 1.012, so the width is set almost entirely by the validation sample size (156 and 267 compounds).}
\end{table}

\begin{table}[h]
\caption{Binning-free metrics give the same ordering as expected calibration error}\label{tab:metrics}
\begin{tabular}{@{}llcccc@{}}
\toprule
Dataset & Control design & ECE & Debiased ECE & ICI & E90 \\
\midrule
DILIrank & Random (standard) & 3.35 & 6.76 & 4.23 & 3.20 \\
 & Random, validation prevalence matched & 1.76 & 2.51 & 1.88 & 1.42 \\
 & Random, both prevalences matched & 1.37 & 1.66 & 1.43 & 1.20 \\
\midrule
DICTrank & Random (standard) & 2.08 & 2.81 & 2.25 & 1.63 \\
 & Random, validation prevalence matched & 1.45 & 1.79 & 1.57 & 1.34 \\
 & Random, both prevalences matched & 1.21 & 1.35 & 1.28 & 1.09 \\
\botrule
\end{tabular}
\footnotetext{Each entry is the temporal-split value divided by the mean of 200 matched random-split controls. ECE uses ten equal-width bins. The debiased ECE subtracts the mean ECE of 1,000 null samples drawn as $y^{*}\sim\mathrm{Bernoulli}(p)$ with the predictions held fixed, which are perfectly calibrated by construction, so the subtracted quantity is the finite-sample bias of ECE. ICI and E90 are the mean and the 90th percentile of $|p-f(p)|$ from a loess calibration curve and use no bins at all. \textbf{All 24 entries exceed one}, and the debiased ratio is the larger of the two in 6 of 6 designs, because the bias is a larger share of the small control error than of the large temporal error. The bias is 14.7\% of the temporal ECE for DILIrank and 15.9\% for DICTrank, and the observed values exceed the 97.5th percentile of their null distributions in both, so the loss is not a small-sample artefact. A Spiegelhalter test rejects calibration for the controls as well ($z=4.31$ to $12.15$): the claim is a difference in degree, not a contrast between calibrated and uncalibrated.}
\end{table}

\begin{table}[h]
\caption{Tuning lowers calibration error but leaves the split-design gap}\label{tab:hparam}
\begin{tabular}{@{}llccccc@{}}
\toprule
Dataset & Configuration & Temporal ECE & Control ECE (R1) & Ratio (R1) & Control ECE (R3) & Ratio (R3) \\
\midrule
DILIrank & Fixed (as used throughout) & 0.4384 & 0.1294 & 3.39 & 0.3247 & 1.35 \\
 & Selected on training era, log loss & 0.4103 & 0.0837 & 4.90 & 0.3128 & 1.31 \\
\midrule
DICTrank & Fixed (as used throughout) & 0.2702 & 0.1272 & 2.12 & 0.2200 & 1.23 \\
 & Selected on training era, log loss & 0.2128 & 0.0888 & 2.40 & 0.1742 & 1.22 \\
\botrule
\end{tabular}
\footnotetext{A grid of 54 configurations crossing tree depth, number of trees, learning rate and minimum child weight was evaluated in full. Configurations were selected on the most recent quarter of the training era, never on the validation set. The fixed configuration used throughout this work ranks 39 of 54 for DILIrank and 39 of 54 for DICTrank by temporal ECE, so the absolute values reported elsewhere are not the best attainable and the grid spans 0.3700--0.4753 and 0.1716--0.3215. \textbf{The ratio to matched random-split controls nevertheless survives tuning} and widens against the standard control, because tuning lowers the control error more than the temporal error. Choosing with hindsight on the validation set itself, 0 and 0 of 54 configurations reach the control level, so the failure is not an artefact of how configurations were chosen. Temporal ECE rises monotonically with depth, number of trees and learning rate in both datasets.}
\end{table}

\begin{table}[h]
\caption{Minimum detectable improvement for each proposed remedy}\label{tab:power}
\begin{tabular}{@{}lllccccc@{}}
\toprule
Dataset & Remedy & Availability & $\Delta$ECE & 95\% CI & MDE & MDE / gap & Verdict \\
\midrule
DILIrank & Uncalibrated (calibration set withheld) & reference & +0.0156 & [-0.027, 0.043] & 0.0494 & 16.1\% & no useful gain \\
 & Platt, training-era calibration set & deployable & -0.0234 & [-0.069, 0.004] & 0.0533 & 17.3\% & no useful gain \\
 & Isotonic, training-era calibration set & deployable & -0.0006 & [-0.051, 0.029] & 0.0567 & 18.4\% & no useful gain \\
 & Prior-shift correction & needs test prevalence & -0.1342 & [-0.179, -0.102] & 0.0559 & 18.2\% & improves \\
 & Recency weighting at selected half-life & deployable & +0.0040 & [-0.025, 0.035] & 0.0427 & 13.9\% & no useful gain \\
 & Class weighting to validation prevalence & needs test prevalence & -0.0880 & [-0.128, -0.051] & 0.0543 & 17.7\% & improves \\
\midrule
DICTrank & Uncalibrated (calibration set withheld) & reference & +0.0233 & [-0.017, 0.045] & 0.0439 & 31.3\% & no useful gain \\
 & Platt, training-era calibration set & deployable & -0.0987 & [-0.137, -0.074] & 0.0463 & 33.0\% & improves \\
 & Isotonic, training-era calibration set & deployable & -0.0898 & [-0.138, -0.067] & 0.0511 & 36.4\% & improves \\
 & Prior-shift correction & needs test prevalence & -0.0273 & [-0.061, -0.012] & 0.0346 & 24.7\% & improves \\
 & Recency weighting at selected half-life & deployable & -0.0444 & [-0.084, 0.002] & 0.0614 & 43.8\% & undecided \\
 & Class weighting to validation prevalence & needs test prevalence & -0.0345 & [-0.059, -0.001] & 0.0405 & 28.9\% & improves \\
\botrule
\end{tabular}
\footnotetext{Each remedy is compared with the uncalibrated model on the same validation set by a paired bootstrap of 4,000 resamples in which the predictions are held fixed, so the two arms remain correlated. The gap a remedy must close is the uncalibrated ECE minus the mean ECE of matched random-split controls (0.3076 for DILIrank, 0.1403 for DICTrank). MDE is $2.80\times\mathrm{SE}(\Delta\mathrm{ECE})$, the smallest improvement that would have reached significance at two-sided $\alpha=0.05$ with 80\% power. \textbf{7 of the 12 arms were powered to detect an improvement of a quarter of the gap}, so the negative findings are not an artefact of sample size; 1 arm remains undecided and it is the one with the largest MDE. Of the 6 remedies that need no knowledge of the validation set, 2 improve calibration significantly, both in DICTrank; none does in DILIrank.}
\end{table}

\begin{table}[h]
\caption{Importance weighting by a domain classifier does not restore calibration}\label{tab:covshift}
\begin{tabular}{@{}llcccccc@{}}
\toprule
Dataset & Training weights & Effective $n$ & Max $w$ & AUROC & ECE & Resolution \\
\midrule
DILIrank & No weighting (as used throughout) & 640.0 & 1.00 & 0.6476 & 0.4384 & 0.0191 \\
 & Density ratio, uncapped & 24.6 & 72.55 & 0.6566 & 0.4755 & 0.0278 \\
 & Density ratio, capped at 20 & 26.4 & 61.27 & 0.6671 & 0.4557 & 0.0156 \\
 & Density ratio, capped at 10 & 42.0 & 36.13 & 0.6553 & 0.4689 & 0.0184 \\
 & Density ratio, capped at 5 & 61.4 & 21.25 & 0.6557 & 0.4670 & 0.0180 \\
 & Density ratio, capped at 2 & 102.8 & 11.31 & 0.6549 & 0.4617 & 0.0144 \\
\midrule
DICTrank & No weighting (as used throughout) & 652.0 & 1.00 & 0.6198 & 0.2702 & 0.0125 \\
 & Density ratio, uncapped & 13.4 & 138.47 & 0.6225 & 0.2904 & 0.0096 \\
 & Density ratio, capped at 20 & 40.9 & 35.77 & 0.6304 & 0.3022 & 0.0131 \\
 & Density ratio, capped at 10 & 61.5 & 21.41 & 0.6374 & 0.2870 & 0.0112 \\
 & Density ratio, capped at 5 & 92.7 & 13.33 & 0.6296 & 0.2912 & 0.0148 \\
 & Density ratio, capped at 2 & 147.3 & 7.19 & 0.6285 & 0.2944 & 0.0100 \\
\botrule
\end{tabular}
\footnotetext{Training compounds are weighted by $w(x)=P(\mathrm{post}|x)/(1-P(\mathrm{post}|x))$, where the probability comes from out-of-fold predictions of a classifier trained to separate the two eras by structure alone (area under the curve 0.6835 and 0.7917). The weighting uses no outcome labels from the validation era, so unlike prior-shift correction it is in principle deployable, provided the structures of the future compounds are available. Effective sample size is $(\sum w)^2/\sum w^2$. \textbf{ECE rises in all 10 of the 10 weighted arms in both datasets}, at every cap, while the effective sample size collapses by up to 97.9\%; the density ratio has a heavy tail and spends the sample. AUROC rises slightly, so the manipulation improves ranking and degrades probabilities. Applying the same weighting to the matched controls as well leaves the ratio at 1.42 and 1.27, against 1.37 and 1.21 unweighted.}
\end{table}

\begin{table}[h]
\caption{At equal shift the temporal split still loses more calibration}\label{tab:shiftmatched}
\begin{tabular}{@{}lllcccccc@{}}
\toprule
Dataset & Split axis & $\alpha^{*}$ & Shift AUC & Target & ECE (matched) & ECE (temporal) & Ratio & SE units \\
\midrule
DILIrank & Molecular weight & 0.29 & 0.6598 & 0.6704 & 0.3295 & 0.4384 & 1.33 & 7.3 \\
 & logP & 0.33 & 0.6753 & 0.6704 & 0.3845 & 0.4384 & 1.14 & 7.2 \\
 & Random direction & 0.31 & 0.6734 & 0.6704 & 0.2809 & 0.4384 & 1.56 & 11.5 \\
\midrule
DICTrank & Molecular weight & 0.46 & 0.7738 & 0.7807 & 0.2189 & 0.2702 & 1.23 & 8.2 \\
 & logP & 0.46 & 0.7672 & 0.7807 & 0.2801 & 0.2702 & 0.96 & -1.8 \\
 & Random direction & 0.46 & 0.7795 & 0.7807 & 0.2005 & 0.2702 & 1.35 & 11.5 \\
\botrule
\end{tabular}
\footnotetext{Non-temporal splits are built by blending the rank along a chemical axis with a random rank, so the amount of shift can be tuned continuously; $\alpha^{*}$ is the blend whose shift matches the temporal split. Shift is measured as the cross-validated area under the curve of a classifier separating the two sides from structure alone, the same quantity reported for the temporal split elsewhere. Every split matches the temporal target to within 0.0134. The size and prevalence of both sides are matched to the temporal split, so the axis is the only remaining difference. \textbf{In 5 of the 6 comparisons the temporal split has the larger ECE}, by 7.2 to 11.5 times the repetition-level standard error; 1 comparison is a tie and none is reversed. The matched splits hold fewer compounds than the temporal validation set, which inflates the small-sample bias of ECE on their side, so the comparison is conservative. Calibration loss under temporal validation therefore does not reduce to the amount of distribution shift.}
\end{table}

\begin{table}[h]
\caption{Annual retraining with recalibration is the one remedy that works}\label{tab:prequential}
\begin{tabular}{@{}llcccccc@{}}
\toprule
Dataset & Design & $n$ & ECE & SD & AUROC & Calibration slope & Mean prediction \\
\midrule
DILIrank & Fixed model (trained once, pre-2010) & 156 & 0.4384 & --- & 0.6476 & 0.2700 & 0.7321 \\
 & Annual retraining & 156 & 0.3751 & --- & 0.6842 & 0.3709 & 0.6677 \\
 & Annual retraining and recalibration & 156 & \textbf{0.2836} & --- & 0.7218 & 0.8705 & 0.5708 \\
 & Shuffled years & 156 & 0.1170 & 0.0216 & 0.7370 & 0.5186 & 0.5921 \\
 & Shuffled years, prevalence matched & 96 & 0.2383 & 0.0231 & 0.7299 & 0.4682 & 0.5218 \\
 & Shuffled years, recalibrated & 156 & 0.0847 & 0.0247 & 0.7167 & 0.9110 & 0.5813 \\
 & Shuffled years, prevalence matched, recalibrated & 96 & 0.2335 & 0.0233 & 0.6999 & 0.7996 & 0.5363 \\
\midrule
DICTrank & Fixed model (trained once, pre-2010) & 267 & 0.2702 & --- & 0.6198 & 0.2223 & 0.8453 \\
 & Annual retraining & 267 & 0.2583 & --- & 0.6227 & 0.2527 & 0.8157 \\
 & Annual retraining and recalibration & 267 & \textbf{0.0960} & --- & 0.5943 & 0.6973 & 0.6985 \\
 & Shuffled years & 267 & 0.1330 & 0.0221 & 0.6879 & 0.3915 & 0.7672 \\
 & Shuffled years, prevalence matched & 191 & 0.1868 & 0.0192 & 0.6916 & 0.4045 & 0.7507 \\
 & Shuffled years, recalibrated & 267 & 0.0538 & 0.0121 & 0.6750 & 0.9103 & 0.7283 \\
 & Shuffled years, prevalence matched, recalibrated & 191 & 0.1158 & 0.0196 & 0.6696 & 0.9089 & 0.7177 \\
\botrule
\end{tabular}
\footnotetext{Each validation year is predicted by a model trained only on compounds approved before it, and the predictions of all years are pooled into a single evaluation set, which is the same set of compounds used throughout this work. Recalibration uses Platt scaling fitted on the most recent quarter of each training window, the same procedure as the static recalibration analysis but on a window that moves. The shuffled-year rows are controls in which the approval years are permuted, so the window sizes are preserved and the time ordering is destroyed; they are reported with and without prevalence matching and with and without recalibration, so that every arm can be read against a control that matches it on both. Paired bootstrap differences against the fixed model are -0.0681 and -0.0200 for retraining alone and -0.1655 and -0.1759 with recalibration (DILIrank and DICTrank datasets); retraining alone is not significant in DICTrank. After retraining and recalibration the error is 1.21 times the matched control for DILIrank and 0.83 times for DICTrank, against 1.37 and 1.21 for the static design. These figures assume that a toxicity label is available in the year of approval, which is optimistic, so they bound what periodic retraining can achieve rather than estimate it.}
\end{table}


%====================================================================%
% Figure environments                                               %
% PDF 파일은 src/make_figures.py 가 figures/ 에 만든다.               %
% 범례에 **반복 횟수나 판정어를 손으로 적지 말 것.** 설정이 바뀌면      %
% 조용히 거짓이 된다 — Fig1 각주의 '200 random partitions' 가 실제로   %
% 그랬다(2026-10-22). 횟수는 Methods 와 본문 표가 결과 파일에서 읽어   %
% 쓴다. 각주에는 설계 상수(구간 10개, 절단 2005/2010/2015, 95%% CI)만   %
% 적는다. 수치를 적어야 하면 verify12 에 등록할 것.                     %
%====================================================================%

\begin{figure}[h]
\centering
\includegraphics[width=0.9\textwidth]{Fig1_split_strategies.pdf}
\caption{Discrimination and calibration under five validation designs}\label{fig:split}
\footnotetext{(a) AUROC under stratified five-fold cross-validation, Bemis--Murcko scaffold five-fold cross-validation and the 2010 temporal split, all on the same compounds. Error bars are the fold-to-fold standard deviation for the cross-validation bars and the 95\% bootstrap interval for the temporal bar. (b) Expected calibration error under the temporal split and the three prevalence-matched random controls R1--R3; error bars are the standard deviation across the prevalence-matched random partitions, whose number is given in Methods.}
\end{figure}

\begin{figure}[h]
\centering
\includegraphics[width=0.9\textwidth]{Fig2_calibration_curves.pdf}
\caption{Reliability curves under random and temporal validation}\label{fig:calib}
\footnotetext{Observed positive fraction against mean predicted probability over ten equal-width bins. The diagonal marks perfect calibration. Points below the diagonal indicate over-confident predictions.}
\end{figure}

\begin{figure}[h]
\centering
\includegraphics[width=0.9\textwidth]{Fig3_robustness.pdf}
\caption{Calibration degradation across models, representations and cut-off years}\label{fig:robust}
\footnotetext{(a) Ratio of temporal to random expected calibration error for all 32 combinations of four learners and four molecular representations, in both datasets. (b) Change in AUROC and ratio of expected calibration error at cut-off years 2005, 2010 and 2015.}
\end{figure}

\begin{figure}[h]
\centering
\includegraphics[width=0.9\textwidth]{Fig4_prevalence_sweep.pdf}
\caption{Calibration error against the held-out positive rate}\label{fig:prev}
\footnotetext{Expected calibration error of random partitions resampled to a range of held-out positive rates, against the temporal value. The temporal point lies above the curve at its own positive rate.}
\end{figure}

\begin{figure}[h]
\centering
\includegraphics[width=0.9\textwidth]{Fig5_recalibration_conformal.pdf}
\caption{What post hoc remedies recover, and where the coverage guarantee breaks}\label{fig:remedy}
\footnotetext{(a) Expected calibration error along the recalibration ladder: no recalibration, Platt scaling and isotonic regression fitted inside the training period, a prior-shift correction, and an oracle Platt fit that uses held-out labels. The oracle is an upper bound, not a deployable method; the shaded band is therefore the most that any in-period remedy could recover. (b) Empirical coverage of split conformal prediction at the two nominal levels, under random and temporal validation, together with coverage on the negative class alone under temporal validation. Dashed lines mark the nominal level.}
\end{figure}

\begin{figure}[h]
\centering
\includegraphics[width=0.9\textwidth]{Fig6_retraining_delay.pdf}
\caption{Periodic retraining with recalibration, and its sensitivity to label delay}\label{fig:retrain}
\footnotetext{(a) Expected calibration error in the expanding-window deployment simulation: a model fitted once before the split, a model refitted every year, and a model refitted and recalibrated every year. (b) Paired change in expected calibration error of the refit-and-recalibrate arm relative to the fixed model, as a function of the assumed number of years a label takes to mature. Shaded bands are 95\% paired bootstrap intervals; values below zero indicate improvement.}
\end{figure}

\begin{figure}[h]
\centering
\includegraphics[width=0.9\textwidth]{FigS1_atc_stratified.pdf}
\caption{Held-out performance stratified by first-level ATC class}\label{fig:atc}
\footnotetext{Classes with at least ten held-out compounds are shown. No subgroup difference remains significant after Benjamini--Hochberg correction within the exploratory family.}
\end{figure}


%%===================================================================%%
%% Discussion and Conclusions — first draft, 2026-09-21               %%
%% Goes between \section{Results} and \section{Methods}.              %%
%% Every number here also appears in results_body_v2.tex; nothing     %%
%% new is computed in this section.                                   %%
%%===================================================================%%

\section{Discussion}

\subsection{What the temporal split costs, and what the cost is made of}

Both benchmarks lose discrimination under a 2010 temporal split: AUROC falls by
0.088 in DILIrank and 0.077 in DICTrank against a matched random reference, with
permutation $p$ of 0.0070 and 0.0050. Numbers of that size are the usual finding
of a time-split study, and on their own they would not warrant a paper. What the
AUROC does not show is that the same split raises expected calibration error from
0.121 to 0.438 and from 0.126 to 0.270, and that the Murphy decomposition assigns
86.1\% and 82.5\% of the increase in Brier score to the reliability term rather
than the resolution term. (The random-split figures quoted here come from a single
matched partition; the 200-partition controls of the next section give 0.131 and
0.130 for the same quantity, and the two differ only in the second decimal place.)

The practical difference between the two failures is that discrimination is a
property of the ranking and calibration is a property of the number. A model that
loses 0.09 of AUROC still ranks most compounds sensibly. A model whose mean
predicted probability is 0.732 on a held-out set with a 0.308 positive rate is
telling a user something false about every compound it scores. In this setting the
number is what gets acted on: a predicted probability of hepatotoxicity is read as
a probability, compared against a threshold, and used to decide whether a compound
advances. Seventy-one of 156 DILIrank held-out compounds carried a predicted
probability above 0.9, at a mean of 0.962, and 40.8\% of them were positive.

The failure is systematic rather than random, which is why the optimal decision
threshold moves to 0.920 and 0.940 rather than staying near 0.5, and why moving
it recovers a substantial part of the classification performance lost at the
default. A threshold that high is not a usable default; the point is the distance
one has to travel to find it, and the fact that nothing in the AUROC signals that
the journey is necessary.

\subsection{Scaffold splitting is not a proxy for temporal validation}

Structural novelty is the usual stand-in for prospective difficulty, on the
reasoning that future compounds are structurally new. In these two datasets the
substitution does not hold.

Measured by AUROC, Bemis--Murcko scaffold cross-validation is not reliably harder
than stratified cross-validation: it raised DILIrank by 0.002 and lowered DICTrank
by 0.022, both within a fold-to-fold standard deviation, and in opposite
directions. Measured by ECE, scaffold and stratified cross-validation agree in
both datasets to within a standard deviation, while the temporal split on the same
compounds gives ECE 3.4 and 2.0 times larger.

The mechanism is visible in the scaffold counts. Of 542 unique scaffolds among the
796 DILIrank compounds, 455 (84\%) appear exactly once; of 638 among the 919
DICTrank compounds, 545 (85\%). Grouping by scaffold cannot separate what is
already separated: when most scaffolds are singletons, a scaffold-grouped fold is
close to a random fold, and the design converges on the thing it was meant to be
harder than. This is a property of drug-like reference sets rather than of these
two datasets, and it should be checked and reported whenever scaffold splitting is
offered as evidence of prospective validity. The check is one line of code and the
number belongs in the Methods of any paper that uses the design.

None of this argues that scaffold splitting is uninformative. It argues that
structural novelty and temporal novelty are different axes, that these datasets
separate them, and that a model can be robust on one while failing on the other.

\subsection{What does explain the collapse, and what does not}

Three candidate explanations were tested directly. One caveat applies throughout:
only the prevalence controls were evaluated against calibration. The chemical-space
and label-maturity analyses measure the discrimination gap alone, so they can rule
an explanation in or out for AUROC but not for ECE.

\textbf{Prevalence shift is a large part of it.} The positive rate falls from
65.5\% to 30.8\% across the DILIrank boundary and from 77.0\% to 61.4\% across the
DICTrank one, and a model calibrated to one base rate is miscalibrated on another
by construction. Matching both base rates in a random split closes 62\% and 66\%
of the ECE excess. Both datasets agree closely on that fraction, which is itself
informative: the effect is not idiosyncratic. But 38\% and 34\% of the excess survive
full prevalence matching, and the residual remains significant after
correction ($q \le 0.0348$). Prevalence is a major term, not the whole term.

Notably, prevalence does not explain any of the \emph{discrimination} loss. Matching
the held-out positive rate to the training rate moved DILIrank AUROC from 0.648 to
0.632, widening the gap rather than closing it. The two failures have different
causes, which is a further reason not to treat AUROC and ECE as interchangeable
summaries of the same phenomenon.

\textbf{Chemical-space novelty is not supported as an explanation of the
discrimination gap.} Held-out
compounds do sit far from the training set --- median nearest-neighbour Tanimoto
0.271, with 73.7\% below 0.4 --- and AUROC rises across the three similarity bins
large enough to evaluate, from 0.543 to 0.711 to 0.857; the highest-similarity bin
holds 12 compounds and was not evaluated. But splitting the held-out set at the
median gives a difference of only
0.068 at a bootstrap $p$ of 0.527, and the correlation between similarity and the
predicted probability of the true class is weak ($\rho = 0.116$, $p = 0.148$).
Neither survives correction. The bin-wise trend is consistent with a real effect
too small to detect at $n = 156$, and it is reported as such rather than as a
finding. As with label maturity, calibration was not recomputed by similarity
stratum, so this analysis says nothing about whether structural distance
contributes to the calibration collapse.

\textbf{Label maturity contributes to the discrimination gap in one dataset and not
the other.} Post-marketing labels accumulate, so young compounds are
under-labelled, and the gradient is present: DILIrank positive rate rises from
14.0\% at seven years or less to 67.7\% at 16--25 years. Restricting the DILIrank
held-out set to compounds with at least twelve years of observation narrows the
AUROC gap by 24\%, but only on 50 compounds, and down-weighting negatives by
observation time changes it by less than 0.003. DICTrank does not move at all: the
most favourable control reduces the gap by 2\%, and the eight-year restriction
widens it. In both datasets the comparison that holds training and held-out
maturity to the same standard --- the only one that is a fair test --- widens the
gap rather than narrowing it.

Two further observations bear on how much weight the maturity story can carry. The
gradient does not continue: in the oldest stratum, which holds the majority of both
datasets, the positive rate falls back, so label accumulation appears to saturate
rather than to continue indefinitely. And the maturity controls were evaluated on
AUROC alone; calibration was not measured under them. Whether label censoring also
contributes to the calibration collapse is therefore untested here, and no claim is
made about it either way. The evidence that prevalence shift is the dominant
calibration term comes from the nested controls of the previous subsection, which
are independent of this one.

\subsection{Recalibration is not a remedy, and the reason matters}

An oracle Platt scaling fitted on the held-out labels reduces ECE to 0.048 and
0.023, so the miscalibration is in principle correctable by a monotone map. Under
a prospective constraint it is not corrected. Internal Platt scaling and isotonic
regression, fitted on the most recent quarter of the training compounds, close
8--10\% of the available gap in DILIrank and 40--42\% in DICTrank.

That the two datasets differ by a factor of four to five is the interesting part.
An internal calibration set fixes the miscalibration a model would have had at the
end of the training period. It cannot anticipate a shift that happens afterwards.
One reading of the difference is that more of the shift falls beyond the last
training year in DILIrank, whose base rate moves by 35 percentage points across the
boundary, than in DICTrank, where it moves by 16. That is a hypothesis consistent
with two datasets, not a tested mechanism, and it would take more endpoints to
distinguish it from other explanations.

A correction that explicitly rescales for a change in base rate recovers 38\% in
DILIrank and 12\% in DICTrank, the reverse ordering, and for the same reason. But
that correction is not prospective: it uses the held-out positive rate, which a
user would have to estimate rather than know. It is reported because it isolates
the base-rate component, not as a method anyone can deploy. Estimating a future
prevalence of drug-induced liver injury among not-yet-approved compounds is not
obviously easier than the prediction problem itself.

The practical conclusion is narrow and, we think, defensible: recalibration on
historical data is worth doing and is not sufficient, and a model card that
reports only internally-calibrated performance is reporting the easy half of the
problem.

\subsection{An assumption-free confirmation}

Split conformal prediction provides a coverage guarantee that depends on
exchangeability and on nothing else --- no distributional assumption, no
calibration metric, no binning choice. A temporal split violates exchangeability
by construction, so the guarantee is expected to fail, and it does: marginal
coverage at a nominal 90\% falls to 73.7\% and 82.4\%, in the same order as the
ECE ratios.

The class-conditional decomposition adds something the marginal figure hides. The
loss falls almost entirely on the negative class --- 63.9\% and 55.3\% against
nominal 90\%, while the positive class stays at or above nominal. In this domain
the negative label is the reassuring one, so a prediction set that omits it is a
model declining to certify as safe a compound the label calls safe. Marginal
coverage reports a guarantee mildly violated; the decomposition shows that one
half of it has gone and the other has not. Split conformal guarantees only
marginal coverage, so these class-conditional numbers are descriptive rather than
tested, which is why the random reference itself can fall below nominal within a
class.

\subsection{Relation to previous work}

Time-split validation has been standard since Sheridan, and has been applied to
DILI prediction before. Calibration degradation under distribution shift has also
been reported: Morger and colleagues documented it across ChEMBL endpoints and
showed that replacing the calibration set mitigates it, and Friesacher and
colleagues showed that temporal shift in industrial pharmaceutical data degrades
uncertainty quantification. Crucitti and Mosquera Orgueira recently argued that the
choice of train--test split matters more for evaluation than for the measured
performance itself, which is the same observation this work makes from a different
direction.

What this study adds is the combination of three things. The endpoints are
regulatory labels for clinical toxicity rather than biochemical activity, and the
time axis is the drug approval year rather than a publication date or a database
release; these are the quantities a regulatory or portfolio decision actually turns
on. The prevalence confound is separated from the rest with nested controls rather
than assumed away. And the calibration claim is made three times over --- through
ECE and the Murphy decomposition, through a distribution-free conformal argument,
and through sensitivity analyses over seeds, binning schemes, cut-off years,
learners and representations --- so that it does not rest on any one metric's
known weaknesses.

\subsection{Limitations}

The held-out sets are small. DILIrank contributes 156 compounds with 48 positives
and a bootstrap AUROC interval of 0.555--0.736. Subgroup cells range from 10 to 204
compounds and none of the twelve tests survives multiple-comparison correction; the
below-chance AUROC in the 32-compound antineoplastic subgroup is reported for
completeness and supports no claim here.

The approval year is a proxy for the date at which a compound entered the world,
not for the date at which its toxicity label was fixed. An audit identified seven
compounds whose earliest matching approval record predates the split boundary, and
three DICTrank compounds recorded as approved after the 2023 label release, which
is more likely an artefact of name matching than a real approval. Applying the seven
corrections moves the ECE ratio by at most 0.31 and the reliability share by at most
2.2 percentage points, so the conclusions do not turn on them. Two of the three
post-2023 compounds are among those seven; the third, selumetinib, is not, because
its earliest matching record is also after the split boundary. The underlying
imprecision in the approval-year proxy remains.

Structure-based deduplication collapses salts and stereoisomers to a common parent,
which merged fourteen groups carrying conflicting labels. Removing them in full is
a defensible rule, fixed before its effect was examined, and its effect is not
one-directional --- re-running the decomposition on the retained pre-filter sets
gives 85.6\% against 86.1\% in DILIrank and 84.5\% against 82.5\% in DICTrank, a
rise in one and a fall in the other --- but it is still an intervention on the
analysis set, and the removed groups are listed in full so a reader can judge it.

ATC assignment carries two unresolved ambiguities: where KEGG lists several codes
for one drug the first alphabetically is taken, and where two drugs normalise to
the same name key the first encountered wins. Both would need the KEGG REST
interface, which was not reachable during this work. The class-L results are
therefore hypothesis-generating regardless of their nominal $p$-values.

The molecular representations are fingerprint- and descriptor-based. Graph neural
networks were not evaluated, on the judgement that a held-out set of 156 compounds
cannot distinguish architectures reliably; that judgement is itself a limitation
rather than a result.

The Murphy decomposition and the conformal analysis were computed for one
model--representation pair. ECE itself was checked across all 32 combinations, so
the calibration claim is not architecture-specific, but the decomposition and the
coverage analysis have not been repeated across the grid.

Two of the three candidate explanations were evaluated on AUROC only. The maturity
controls establish that censoring explains at most a quarter of the discrimination
gap in DILIrank and essentially none of it in DICTrank, and the similarity analysis
finds no significant relation between structural distance and per-compound accuracy,
but neither recomputes calibration. The claim that prevalence shift is the dominant
calibration term therefore rests on the nested prevalence controls alone. Repeating
the maturity and similarity analyses with ECE as the outcome would test it from a
second direction and is the most obvious next experiment.

Finally, the three-family structure used for multiple-comparison correction was
fixed after the analyses were run rather than pre-registered. The grouping criteria
are stated explicitly so that a reader can regroup the tests and recompute; the
confirmatory family survives at $q = 0.0070$ and the exploratory family yields
nothing significant under any grouping we examined.

\subsection{Implications for practice}

Three recommendations follow, in decreasing order of confidence.

Report calibration alongside discrimination whenever a temporal or prospective
split is used. In these data ECE was worse under the temporal split in all 32
model--representation combinations and at all three cut-off years, with the ratio
never falling below 1.6 across cut-offs and below 1.43 across the grid. Over the
same conditions $\Delta$AUROC varied by a factor of 2.4 across cut-offs in DILIrank,
swapped the rank of its 2005 and 2015 values between the two datasets, and changed
sign in one of the 32 combinations. The calibration signal was the more consistent of the two, and it
is the one usually omitted.

Do not treat scaffold splitting as evidence of prospective validity without
reporting the scaffold-size distribution. If most scaffolds are singletons, the
design is closer to a random split than its name suggests.

Treat an internally calibrated model as calibrated only for the period it was
calibrated on. Where a prospective use is intended, the relevant question is not
whether the probabilities were correct at the end of the training window but
whether anything is known about how the base rate will move afterwards.

\section{Conclusions}

Across two FDA regulatory toxicity benchmarks, splitting on approval year degrades
model performance in a way that AUROC substantially understates. The loss in
discrimination is 0.088 and 0.077 of AUROC; against a random split of the same size
expected calibration error rises from 0.131 to 0.438 and from 0.130 to 0.270,
ratios of 3.35 and 2.08, and the Murphy decomposition attributes 86.1\% and 82.5\%
of the increase in Brier score to the calibration term. The pattern held in all 32 combinations of four learners and four molecular
representations, across three cut-off years, ten random seeds and eight binning
schemes, and was confirmed independently by a distribution-free conformal coverage
analysis that requires no calibration metric at all.

Prevalence shift accounts for 62\% of the calibration collapse in DILIrank and 66\%
in DICTrank, and not for the rest. Chemical-space novelty and label censoring were tested against the discrimination
gap only. Chemical-space novelty accounts for none of it; label censoring accounts
for at most a quarter of it in DILIrank and essentially none in DICTrank. Whether
either contributes to the calibration collapse was not measured.
Bemis--Murcko scaffold splitting does not reproduce it, because in drug-like
reference sets most scaffolds occur once and the design reduces to something close
to a random split.

The practical consequence is that a model reported as losing "only" 0.09 of AUROC
under time-split validation may be producing probabilities that are wrong by a
factor that no reported number reveals. Calibration metrics are cheap to compute
and, in this setting, carry the more stable and more actionable signal. They should
be reported as a matter of course.


% ======================================================================
% Methods -- Journal of Cheminformatics submission
%
% HOW TO USE
%   1. Open the official Springer Nature LaTeX template on Overleaf.
%   2. Delete the template's placeholder body text.
%   3. Paste everything below this header into sn-article.tex, or save this
%      file as methods.tex and add \input{methods} there.
%
% CITATIONS
%   Numbers in square brackets follow the finalised 1-72 reference list.
%   To move to BibTeX later, replace each [n] with \cite{key}.
%
% All characters below are plain ASCII; compiles with pdflatex.
% ======================================================================

\section{Methods}

\subsection{Datasets and curation}

Hepatotoxicity labels came from DILIrank 2.0 [19, 20], the 2025 revision of the original DILIrank list [18], and cardiotoxicity labels from DICTrank [14, 21].

Both sources grade concern on an ordinal scale and both carry an ambiguous grade. I treated the two upper grades as positive and the lowest as negative, and discarded the ambiguous compounds: they carry no directional information, and forcing them into either class would put unverifiable labels into the evaluation set. For DILIrank 2.0 this removed 354 of 1,336 entries, leaving 982 compounds (568 positive, 414 negative). DICTrank lists 1,318 entries, of which 107 are ambiguous and a further 27 carry no active-ingredient name; 1,184 remained (841 positive, 343 negative).

Structures came from PubChem [24], resolved by compound name. Each retrieved SMILES was standardised in RDKit [27] --- split the molecule into fragments, keep the largest fragment containing carbon, re-canonicalise --- which strips counter-ions and salt forms. Compounds that could not be resolved or parsed were dropped. Duplicates were then removed on the first 14 characters of the InChIKey, collapsing salt forms, hydrates and stereoisomers of a single parent structure into one record. Of the 982 labelled DILIrank compounds, 891 resolved to a structure and 870 survived deduplication; for DICTrank the corresponding figures are 1,087 and 968.

Approval years came from the Drugs@FDA bulk files [28]. \texttt{Products.txt} was joined to \texttt{Submissions.txt} on application number, keeping only approved submissions with a parsable action date and taking the earliest such date per application; combination products were split on the semicolon separator so that each active ingredient carried its own date. Matching drug names to ingredient names is the fragile step. I used a cascade that begins strict and loosens: exact match after normalisation, then match after stripping trailing salt and hydrate tokens, then match on the first token, and --- for DILIrank only --- substring containment restricted to names of at least five characters. Of the 870 structurally curated DILIrank compounds, 795 matched exactly, 11 through the looser stages and 64 not at all; I checked all 11 loose matches against the drugs' known approval histories. The cascade differs slightly between the two sources because their ingredient strings are formatted differently, but this has no consequence here: applying the DILIrank version to DICTrank left every compound's year unchanged and moved none across the split boundary.

This left \textbf{796 DILIrank compounds} (approved 1939--2022) and \textbf{919 DICTrank compounds} (1939--2025).

\subsection{Molecular representation}

Each compound was represented by a 2,048-bit Morgan fingerprint of radius 2 [29] concatenated with six RDKit descriptors: molecular weight, logP, topological polar surface area, hydrogen-bond donors, hydrogen-bond acceptors and rotatable bonds [27]. The robustness grid additionally uses the fingerprint alone, the 167-bit MACCS keys [30], and the six descriptors alone. The descriptor block is identical wherever it appears, so the fingerprint-plus-descriptor cell of that grid reproduces the primary analysis exactly.

\subsection{Models}

The primary model is XGBoost [31]: 400 trees, depth 6, learning rate 0.05, subsample and column-subsample both 0.8, logistic loss, seed 42. The robustness grid adds a random forest [32] (500 trees, at least two samples per leaf), L2-penalised logistic regression (C = 1, up to 3,000 iterations, features scaled but not centred so the fingerprint stays sparse) and LightGBM [33] (400 trees, 31 leaves, learning rate 0.05, subsample and column-subsample 0.8). All were fitted through the scikit-learn interface [34].

I did not tune hyperparameters. One fixed configuration runs everywhere, so that a difference between validation schemes is a property of the scheme rather than of how hard the search was under each one. Nor was any class weighting, resampling or threshold adjustment applied at fit time. The predicted probability is the object of study here, and correcting imbalance during fitting would alter precisely the quantity under examination.

\subsection{Validation designs}

\paragraph{Temporal split.}
Compounds approved before 2010 form the training set, those approved in 2010 or later the held-out set: 640 and 156 for DILIrank (positive rate 65.5\% training, 30.8\% held out), 652 and 267 for DICTrank (77.0\% and 61.4\%).

\paragraph{Random reference.}
A single random partition is noisy, so the random reference is an average over repeated partitions drawn at the same held-out size as the temporal split. The number of repeats varies by analysis --- 20 for the calibration analysis, the Brier decomposition and the conformal analysis; 200 for the prevalence-matched controls; 15 for the model-by-representation grid; 1000 for the permutation test --- and is stated with each result. The resulting estimates are close but not identical: for DILIrank the random reference is 0.735 over the 1000 permutations of the significance test and 0.738 over the 200 repeats of the prevalence-matched controls. Cross-validation, used only in the split-strategy comparison below, is a different procedure, and is called cross-validation throughout.

\paragraph{Split strategies.}
To separate chemical novelty from temporal shift I compared five-fold stratified cross-validation against five-fold Bemis--Murcko scaffold cross-validation [58, 59], run on the same compounds as the temporal split. Scaffolds were taken without stereochemistry and assigned greedily, largest group first, to whichever fold was smallest at the time, so no scaffold spans folds. All three strategies pass through the same features, model and metrics.

\paragraph{Prevalence-matched controls.}
The temporal split changes the positive rate between training and held-out data --- prior probability shift in the taxonomy of Moreno-Torres et al. [62] --- which confounds any calibration comparison against an ordinary random split. I therefore built three nested controls, each averaged over 30 partitions. R1 is an unconstrained random split. R2 subsamples the held-out set to the positive rate of the temporal held-out set. R3 subsamples both training and held-out sets to the corresponding temporal rates. Subsampling keeps every instance of the limiting class and draws the other without replacement, at the largest size that reaches the target rate. Partitions that fell below 30 held-out or 100 training compounds, or that lost a class entirely, were discarded.

\subsection{Evaluation metrics}

Discrimination is reported as AUROC [63], balanced accuracy and the Matthews correlation coefficient [60, 61], the latter two at the default threshold of 0.5 and again at the threshold maximising Youden's J.

Probability quality is reported as the Brier score [35]; expected and maximum calibration error over ten equal-width bins [38--40]; and the calibration slope and intercept obtained by regressing the outcome on the predicted log-odds, where a slope of 1 and an intercept of 0 mark perfect calibration [41].

\paragraph{Brier decomposition.}
Which of the two carries the loss is easy to assert and harder to show, so I decomposed the Brier score in Murphy's form [36, 37] over the same ten bins:

\begin{equation}
\mathrm{BS} = \mathrm{reliability} - \mathrm{resolution} + \mathrm{uncertainty}.
\end{equation}

Uncertainty is fixed by the held-out class balance and is not a property of the model. Resolution rewards separating compounds and tracks discrimination; reliability penalises predicted probabilities that do not match observed frequencies and tracks calibration. The change in Brier score between the random reference and the temporal split is attributed to the two model-dependent terms, and the reliability share reported as $\Delta$reliability / ($\Delta$reliability + $\Delta$resolution). Bin-wise estimates are unstable at these sample sizes, so the random reference for this analysis is pooled across 20 partitions before binning.

\subsection{Recalibration}

If the loss is concentrated in calibration, the natural question is whether it can be repaired --- and repaired without the held-out labels, which in a prospective setting do not exist yet. I split the pre-2010 training data by approval year and held out the most recent 25\% as an internal calibration set, so that the calibration set sits in the future relative to the fitting set while remaining entirely within the training era.

Five strategies were compared: (A) no recalibration; (A$'$) no recalibration, refitted without the calibration set, as a like-for-like baseline; (B) Platt scaling [42] fitted on the internal calibration set; (C) isotonic regression [43, 44] on the same set; (D) a prior-probability correction that shifts the predicted log-odds by the log ratio of target to source odds, the intercept-only form of the Saerens procedure [45, 46]; and (E) an oracle Platt scaling fitted directly on the held-out labels. (E) is not available in practice and is reported only as a ceiling. Recovery is expressed as the fraction of the gap between the uncorrected and oracle ECE that a strategy closes.

\subsection{Conformal prediction}

Parametric recalibration might fail for reasons peculiar to the parametric form, so I repeated the question in a distribution-free setting. Split conformal prediction [47, 48, 53, 54] was applied using the same fit and calibration partition as above, with non-conformity $1 - \hat{p}$ for the true class and the conformal quantile taken at $\lceil (n+1)(1-\alpha) \rceil / n$ of the calibration scores; a prediction set contains every class whose non-conformity falls at or below that quantile. Marginal coverage, class-conditional coverage, mean set size and the empty-set rate were evaluated at target coverages of 90\% and 80\% ($\alpha = 0.10$ and $0.20$), against a random reference averaged over 20 partitions with identical fit, calibration and held-out sizes.

Class-conditional coverage is reported separately because marginal coverage hides which class absorbs the shortfall. Conformal procedures adapted to covariate and label shift do exist [51, 52], but they require the shift to be characterised in advance. The exchangeability assumption behind the coverage guarantee is precisely what a temporal split breaks [49, 50], so a failure of the standard procedure is expected in direction; the question is its size and its asymmetry.

\subsection{Label maturity}

Post-marketing safety labels accumulate over time, so a recently approved drug is more likely to carry a negative label simply because nothing has been reported yet [72]. This is a plausible competing explanation for temporal degradation, and I tested it directly. Observation time was defined as 2025 minus approval year, 2025 being the release year of DILIrank 2.0.

Three analyses follow. First, the held-out set was restricted to compounds with at least 8, 10, 12 or 15 years of observation and the gap recomputed. Second, negative labels were down-weighted in proportion to observation time --- weight $\min(\mathrm{exposure}/\tau,\,1)$ for $\tau$ of 5, 8, 10 and 15 years, with positives always weighted 1 --- and a pairwise-weighted AUROC computed from them. Third, the same observation-time threshold was applied to training and held-out sets together, so that the comparison is not distorted by restricting only one side. Restrictions leaving fewer than 30 held-out compounds, or a single class, were not evaluated.

\subsection{Robustness and sensitivity}

I varied four things. \textbf{Model and representation:} all 16 combinations of the four learners and four representations, for each dataset, at the 2010 cut-off. \textbf{Cut-off year:} the temporal split repeated at 2005, 2010 and 2015, with the Brier decomposition recomputed at each. \textbf{Seed:} the primary pipeline rerun under ten seeds (42, 0, 1, 7, 13, 21, 100, 123, 777, 2024), recording AUROC, ECE, reliability and resolution. \textbf{Binning:} ECE and the decomposition recomputed over 5, 10, 15 and 20 bins under both equal-width and equal-frequency binning, since ECE is known to move with the binning scheme [40].

\subsection{Subgroup analyses}

Held-out performance was stratified by first-level ATC class [26] and by four physicochemical cut-points (molecular weight above 500, four or more aromatic rings, six or more nitrogens, logP of 3 or above). ATC assignments came from a bulk pull of the KEGG DRUG--ATC linkage [25], matched to compound names by exact match after normalisation and then by match after salt-token stripping; the assignment was verified against compounds of known class before use. This matching is deliberately more conservative than the approval-year cascade, since a misassigned therapeutic class would contaminate a whole subgroup rather than move one compound across a split. Chemical-space displacement was quantified as the maximum Tanimoto similarity [22, 23] of each held-out compound to the training set, computed on the same fingerprints.

These analyses are exploratory. They are reported with multiplicity correction and are not used to support any primary claim.

\subsection{Statistical analysis}

The primary test of temporal degradation is a permutation test [55]: 200 random partitions of the same held-out size were drawn, the model refitted on each, and the observed temporal AUROC referred to the resulting null distribution as a one-sided p-value. Confidence intervals for AUROC and ECE are percentile bootstrap (1,000 resamples for AUROC in the primary split, 2,000 for ECE and for subgroup differences) [57]. Differences between independent held-out subgroups --- high versus low similarity, ATC class L versus the rest --- were tested by bootstrapping each subgroup separately and taking the two-sided proportion of sign-consistent resamples. The relationship between training-set similarity and the predicted probability of the true class was tested by Spearman correlation.

I controlled multiplicity with the Benjamini--Hochberg procedure [56], applying it within three families rather than across all tests at once, because the tests are not of one kind. Family 1 holds the two primary confirmatory tests of temporal degradation, one per dataset. Family 2 holds the six prevalence-matched comparisons. Family 3 holds the eight exploratory subgroup tests, which is where the multiplicity problem actually lives. Tests reported at $p < 10^{-4}$ sit at the resolution limit of the resampling and were entered as $10^{-4}$.

\subsection{Software and data availability}

Everything was run under Python 3.11.14 with RDKit 2026.03.5 [27], scikit-learn 1.9.0 [34], XGBoost 3.2.0 [31], LightGBM 4.7.0 [33], SciPy 1.17.1, statsmodels 0.15.0, pandas 3.0.5 and NumPy 2.4.6; figures were drawn in Matplotlib at 600 dpi. All random number generation is seeded, and I re-executed the robustness pipeline independently to confirm that the reported values reproduce exactly.

The datasets are public [20, 21, 28] and the analysis code is deposited at [REPOSITORY DOI --- to be inserted].


%====================================================================%
% Back matter — abbreviations and declarations                       %
%                                                                    %
% 항목과 순서는 저널 지침을 따른다 (2026-10-20 확인):                  %
%   Availability of data and materials / Competing interests /        %
%   Funding / Authors' contributions / Acknowledgements /             %
%   Authors' information(선택)                                        %
% 해당 없는 항목은 제목을 남기고 "Not applicable" 이라고 적는다.        %
% 이 저널은 Ethics approval·Consent for publication 을 요구하지 않는다. %
%====================================================================%

\backmatter

\section*{List of abbreviations}

AUROC, area under the receiver operating characteristic curve; BH,
Benjamini--Hochberg; CI, confidence interval; DICT, drug-induced
cardiotoxicity; DILI, drug-induced liver injury; ECE, expected calibration
error; MCC, Matthews correlation coefficient; MCE, maximum calibration error.
% [최종 원고가 쓰는 약어에 맞춰 손질할 것.]

\section*{Declarations}

\bmhead{Availability of data and materials}

% [TODO — GitHub 저장소와 Zenodo DOI 가 확정되면 채운다.]
The datasets supporting the conclusions of this article, and the analysis code
required to reproduce every result reported here, are available at
[repository URL] and archived at [Zenodo DOI]. The source data are the public
FDA DILIrank and DICTrank releases; the curation pipeline that reproduces the
analysis sets from them is included in the repository.

\bmhead{Competing interests}

The author declares no competing interests.

\bmhead{Funding}

% [TODO — 연구비가 없으면 아래 문장을 그대로 둔다.]
No funding was received for this work.

\bmhead{Authors' contributions}

The author designed the study, implemented and performed all analyses, and
wrote the manuscript.

\bmhead{Acknowledgements}

% [TODO — 대규모 언어모델을 원고 준비에 썼다면 Springer Nature 정책상
% 여기에 도구 이름과 용도를 밝혀야 한다. 쓰지 않았다면 "Not applicable".]
Not applicable.

% \bmhead{Authors' information}  % 선택 — 필요 없으면 지운다


\bibliography{sn-bibliography}

\end{document}
